\documentclass[letterpaper, 10 pt, conference]{ieeeconf}
\IEEEoverridecommandlockouts                          
\overrideIEEEmargins                                      % Needed to meet printer requiremen

% The preceding line is only needed to identify funding in the first footnote. If that is unneeded, please comment it out.
\let\labelindent\relax
\usepackage[inline]{enumitem}
\usepackage{amsmath}
\usepackage{cite}
\usepackage{amssymb,amsfonts}
\usepackage[ruled,vlined]{algorithm2e}
\usepackage{graphicx}
\usepackage{subfloat}
\usepackage{textcomp}
\usepackage{xcolor}
% \usepackage{algorithm}
\usepackage{verbatim} 
\usepackage{mathrsfs}
\usepackage{subfigure}
\usepackage{empheq}
\usepackage{courier}
\usepackage{lipsum}
\usepackage{comment}
\usepackage{mathtools, cuted}
\usepackage{bm}
\usepackage{ wasysym } % For /ocircle
\usepackage{soul}
\usepackage[colorlinks,citecolor=red,urlcolor=black,bookmarks=false,hypertexnames=true,hidelinks]{hyperref}

\def\BibTeX{{\rm B\kern-.05em{\sc i\kern-.025em b}\kern-.08em
    T\kern-.1667em\lower.7ex\hbox{E}\kern-.125emX}}

% \usepackage{tikz,xcolor,hyperref}
\newcommand{\qm}[1]{{\textcolor{cyan}{\textbf{[Qian]} #1}}}
\newcommand{\hkc}[1]{{\textcolor{blue}{[HKG]\emph{\bf#1 }}}}
\newcommand{\todo}{\textcolor{red}{TODO }}
\long\def\revised#1{{\color{blue}#1}}

% Caligraphic

\newcommand{\call}{\mathcal}

\def\cA{{\call A}}
\def\cB{{\call B}}
\def\cC{{\call C}}
\def\cD{{\call D}}
\def\cE{{\call E}}
\def\cF{{\call F}}
\def\cH{{\call H}}
\def\cI{{\call I}}
\def\cJ{{\call J}}
\def\cK{{\call K}}
\def\cL{{\call L}}
\def\cM{{\call M}}
\def\cN{{\call N}}
\def\cO{{\call O}}
\def\cP{{\call P}}
\def\cQ{{\call Q}}
\def\cR{{\call R}}
\def\cS{{\call S}}
\def\cT{{\call T}}
\def\cU{{\call U}}
\def\cV{{\call V}}
\def\cW{{\call W}}
\def\cX{{\call X}}
\def\cY{{\call Y}}
\def\cZ{{\call Z}}

\newcommand{\integers}{\mathbb{Z}}
\newcommand{\naturals}{\mathbb{N}}

\newcommand{\ap}{AP}
\newcommand{\true}{{\tt{True}}}
\newcommand{\false}{{\tt{False}}}
\newcommand{\var}{\pi}
\newcommand{\inp}{\cI}
\newcommand{\inpsyms}{\inp_\textrm{c}}
\newcommand{\inpuser}{\inp_\textrm{u}}
\newcommand{\inpuserprov}{\inp_\textrm{user}}
\newcommand{\inpcen}{\inp_{\textrm{robot}}}
\newcommand{\inpcentrans}{\inp_{\textrm{robot}}^{\textrm{trans}}}
\newcommand{\inpcenorien}{\inp_{\textrm{robot}}^{\textrm{orien}}}
\newcommand{\inpterrain}{\inp_{\textrm{terrain}}}
\newcommand{\inpterrainbool}{\inp_{\textrm{terrain}}^{\cB}}
\newcommand{\inpterrainint}{\cE_\naturals}
\newcommand{\terrainexprs}{E}
\newcommand{\inpstatesterrainexpr}{2^{\inpterrain}}
\newcommand{\inpstatesterrainexprposs}{\dom_{\textrm{terrain}}^{\textrm{poss}}}

\newcommand{\inpstateterrainexpr}{\inpstatesterrain}
\newcommand{\inprequest}{\inp_{\textrm{req}}}
\newcommand{\inprequesttrans}{\inp_{\textrm{req}}^{\textrm{trans}}}
\newcommand{\inprequestorien}{\inp_{\textrm{req}}^{\textrm{orien}}}
\newcommand{\tvar}[2]{n_{#1}^{#2}}
\newcommand{\cells}{\cC}
\newcommand{\xcells}{\cX}
\newcommand{\ycells}{\cY}
\newcommand{\out}{\cO}
\newcommand{\skillcustom}[1]{\skill_{#1}}
\newcommand{\skillprecustom}[1]{\Sigma_{#1}^{\textrm{pre}}}
\newcommand{\skillpostcustom}[1]{\Sigma_{#1}^{\textrm{post}}}
\newcommand{\skill}{o}
\newcommand{\skillpre}{\skillprecustom{\skill}}
\newcommand{\skillpost}{\skillpostcustom{\skill}}
\newcommand{\skillpretrans}{\skillprecustom{\skill_{\textrm{trans}}}}
\newcommand{\skillposttrans}{\skillpostcustom{\skill_{\textrm{trans}}}}
\newcommand{\skillpreorien}{\skillprecustom{\skill_{\textrm{orien}}}}
\newcommand{\skillpostorien}{\skillpostcustom{\skill_{\textrm{orien}}}}

\newcommand{\dom}{\Sigma}
\newcommand{\statevar}{\sigma}
\newcommand{\inpstate}{\statevar_\inp}
\newcommand{\inpstates}{\dom_\inp}
\newcommand{\inpstateterrain}{\statevar_{\textrm{terrain}}}
\newcommand{\inpstateterrainbool}{\statevar_{\textrm{terrain}}^{\cB}}

\newcommand{\inpstateterrainint}{\statevar_{\inpterrainint}}
\newcommand{\inpstatesterrain}{\dom_{\textrm{terrain}}}
\newcommand{\inpstatesterrainint}{\dom_{\inpterrainint}}
\newcommand{\inpstatecen}{\statevar_{\textrm{robot}}}
\newcommand{\inpstatescen}{\dom_{\textrm{robot}}}
\newcommand{\inpstaterequest}{\statevar_{\textrm{req}}}
\newcommand{\inpstatesrequest}{\dom_{\textrm{req}}}
\newcommand{\inpstatesrequestposs}{\dom_{\textrm{req}}^{\textrm{poss}}}
\newcommand{\inpstateterrainposs} {\Sigma_{\inpterrain}^{\textrm{poss}}}

%%%% Ref %%%%
\newcounter{examplecounter}
\setcounter{examplecounter}{0}

\newcommand{\examplelabel}[1]{%
  \refstepcounter{examplecounter}%
  \label{#1}%
  \theexamplecounter%
}

\newcounter{problemcounter}
\setcounter{problemcounter}{0}

\newcommand{\problemlabel}[1]{%
  \refstepcounter{problemcounter}%
  \label{#1}%
  \theproblemcounter%
}

\newcounter{defcounter}
\setcounter{defcounter}{0}

\newcommand{\deflabel}[1]{
    \refstepcounter{defcounter}%
    \label{#1}%
    \thedefcounter%
}

%%%%%%%%%%%%%%%%%%%%%%%%%%
%%%% Spec related %%%%
\newcommand{\envspec}{\varphi_\textrm{e}}
\newcommand{\boolspec}{\varphi_{b}}
\newcommand{\sysspec}{\varphi_\textrm{s}}
\newcommand{\envinit}{\varphi_\textrm{e}^\textrm{i}}
\newcommand{\sysinit}{\varphi_\textrm{s}^\textrm{i}}
\newcommand{\alphainit}{\varphi_\alpha^\textrm{i}}
\newcommand{\envsafety}{\varphi_\textrm{e}^\textrm{t}}
\newcommand{\envterrain}{\varphi_\textrm{e}^\textrm{terrain}}
\newcommand{\alphasafety}{\varphi_\alpha^\textrm{t}}
\newcommand{\envsafetyvio}{\varphi_\textrm{violated}}
\newcommand{\envsafetynothard}{\varphi_\textrm{e}^\textrm{t,skill}}
\newcommand{\envsafetyhard}{\varphi_\textrm{e}^\textrm{t,hard}}
\newcommand{\alphasafetynothard}{\varphi_\alpha^\textrm{t,skill}}
\newcommand{\alphasafetyhard}{\varphi_\alpha^\textrm{t,hard}}
\newcommand{\envsafetyharduser}{\varphi_\textrm{e}^\textrm{t,user}}
\newcommand{\syssafety}{\varphi_\textrm{s}^\textrm{t}}
\newcommand{\syssafetynothard}{\varphi_\textrm{s}^\textrm{t,skill}}
\newcommand{\syssafetyhard}{\varphi_\textrm{s}^\textrm{t,hard}}
\newcommand{\syssafetyuser}{\varphi_\textrm{s}^\textrm{t,user}}
\newcommand{\syssafetycontinue}{\varphi_\textrm{s}^\textrm{t,continue}}
\newcommand{\syssafetychoose}{\varphi_\textrm{s}^\textrm{t,choose}}
\newcommand{\envlive}{\varphi_\textrm{e}^\textrm{g}}
\newcommand{\syslive}{\varphi_\textrm{s}^\textrm{g}}
\newcommand{\alphalive}{\varphi_\alpha^\textrm{g}}
\newcommand{\changecons}{\varphi_\textrm{poss-changes}}
\newcommand{\notallowedrepair}{\varphi_\textrm{disallow}}
\newcommand{\disallowedtrans}{\notallowedrepair}
\newcommand{\specsafe}{\spec^\textrm{t}}
\newcommand{\speclive}{\spec^\textrm{g}}
\newcommand{\alphasafetybool}[1]{\specsafe_{\textrm{a}, #1}}
\newcommand{\alphalivebool}[1]{\speclive_{\textrm{a}, #1}}

\newcommand{\spec}{\varphi}
\newcommand{\spectask}{\spec_{\textrm{task}}}
\newcommand{\specrelax}{\spec_{\textrm{relax}}}
\newcommand{\specrepair}{\spec_{\textrm{repair}}}
\newcommand{\oldspec}{\varphi_{\textrm{old}}}
\newcommand{\newspec}{\varphi_{\textrm{new}}}
\newcommand{\psicheck}{\psi_i^{\textrm{check}}}
\newcommand{\envrelaxhard}{\psi_{\textrm{r}}^{\textrm{hard}}}
\newcommand{\envrelaxnothard}{\psi_{\textrm{r}}^{\textrm{skill}}}
\newcommand{\trace}{\tau}
\newcommand{\traceinf}{\tau_{\textrm{inf}}}
\newcommand{\tracefin}{\tau_{\textrm{fin}}}


%%%%%%%%%%%%%%%%%%%%%%%%%%%%%%%%%

%%%% Functions %%%%
\newcommand{\downright}[3]{\textit{#1}_{#2}(#3)}
\newcommand{\form}[2]{\downright{form}{#1}{#2}}
\newcommand{\formplain}{\textit{form}}
\newcommand{\range}[1]{[#1]}

%%%% Repair %%%%
\newcommand{\winning}{Z}
\newcommand{\strategy}{\cA}

\newcommand{\grounding}{\textrm{G}}
% \newcommand{\groundingreg}{\textrm{G}_\reg}
\newcommand{\inversegrounding}{\grounding^{-1}}
\newcommand{\ps}{\cZ}
\newcommand{\physicalstate}{z}
\newcommand{\reals}{\mathbb{R}}

\makeatother
\begin{document}

\newcommand{\vg}[1]{\bm{#1}}
\renewcommand{\v}[1]{\mathbf{#1}}

% \title{Reactive Gait Synthesis and Footstep Planning for Safety-Critical Quadrupedal Locomotion\\
\title{\LARGE \bf Physically-Feasible Reactive Synthesis for Terrain-Adaptive Locomotion via Trajectory Optimization and Symbolic Repair\\
}

\author{
\large Ziyi Zhou$^{1}$,
Qian Meng$^{2}$, 
% Eohan George$^{3}$,
Hadas Kress-Gazit$^{2}$, and
Ye Zhao$^{1}$
% <-this % stops a space
\thanks{$^{1}$Ziyi Zhou, and Ye Zhao are with Georgia Institute of Technology. {\tt\small \{zhouziyi, yzhao301\}@gatech.edu.}}
\thanks{$^{2}$Qian Meng and Hadas Kress-Gazit are with Cornell University. {\tt\small \{qm34,hadaskg\}@cornell.edu.}}
% \thanks{$^{3}$Eohan George is Principal Investigator at SkyMul Inc. {\tt\small eg@skymul.com.}}
}
\maketitle

\begin{abstract}
% Challenging but safety-critical scenarios, like hazardous debris or scattered construction sites, demand that robots respond to dynamic environments not only versatilely, but also safely and promptly. Although legged robots excel on rough terrains with proven versatility, a formal guarantee of their safety and reactivity remains a significant challenge, hindering their deployment in urgent applications.

% In this work, we adopt linear temporal logic (LTL)-based reactive synthesis to abstract the terrain-adaptive dynamic locomotion problem as a symbolic level gait design and decision-making problem.
This study proposes an integrated planning framework for quadrupedal locomotion over dynamically changing, unforeseen terrains.
Existing terrain-adaptive locomotion approaches either rely heavily on heuristics for reactive foothold selection, compromising safety and reliability, or directly solve a holistic optimization, suffering from expensive computation as terrain features become complex.
% \qm{maybe we could fill this up with heuristic methods and their disadvantages}, 
% To address these challenges, our approach integrates reactive synthesis and Mixed-Integer Convex Program (MICP) to leverage the strength of both approaches. 
To address these challenges,
% Specifically, 
we leverage reactive synthesis to 
generate correct-by-construction controllers
to navigate the robot with proper locomotion gaits at the symbolic level, 
and MICP to ensure dynamic footstep planning with physical feasibility for each symbolic transition.
% \qm{We should add a few sentences to summarize our approach: a global planner to produce a nominal trajectory. MICP,  symbolic repair to guide the search of gait configuration, and online execution that can generalize to unseen terrain combinations.}
% while leveraging convex optimization to provide global feasibility certificate at the motion planning level. 
% The mismatch between the robot's symbolic-level decisions and its physical-level capabilities in navigating adversarial environments is addressed in both the offline and online phases.
% To address the mismatch between the robot's symbolic-level decisions and its physical-level capabilities in maneuvering through adversarial environments, we propose an offline-online framework. 
% A Mixed-Integer Convex Program (MICP) generates motion primitves and offers an approximate dynamics-feasibility guarantee based on terrain information given a prior. 
To handle unrealizable specifications due to limited feasible symbolic transitions, we leverage a 
% closed-loop 
symbolic repair process to identify dynamically feasible missing skills in both offline and online phases. %, thereby ensuring that the synthesis remains realizable.
During the online execution phase, we re-run the MICP using real-world terrain data to bridge the gap between offline synthesis and online execution, along with online symbolic repair to handle unforeseen terrains.
Extensive simulation and hardware experiments demonstrate the efficacy of our framework in discovering missing locomotion skills and enabling prompt reactions to safety-critical environments, such as scattered stepping stones and rebars.
\end{abstract}

\section{Introduction}
% New outline:

% 1. The primary motivation for using reactive synthesis in terrain-adaptive locomotion lies in its contributions to safety and computational efficiency. First, reactive synthesis provides a formal method to ensure safety by systematically considering possible interactions between the robot and surrounding adversarial terrains. Second, it abstracts complex and rich features of the environment, enabling faster computation compared to traditional methods that directly reason about the intricate relationship between the environment and robot dynamics.

% 2. Symbolic repair is crucial for guiding the search for necessary gait configurations traversing adversial terrains. One reason is that it reduces the effort needed to identify all essential skills required for different terrains. Additionally, symbolic repair facilitates the accumulation of experience. As repairs are conducted, the robot identifies and stores necessary skills, progressively enhancing its capabilities as it encounters a wider variety of terrains.

% 3. During the online phase, we address the gap between discrete abstraction and real-world information for both the robot and its environment. This is achieved by solving the MICP in real-time using up-to-date robot and terrain data. Online repair also utilizes real-time terrain information to handle previously unseen environments that were not encountered during offline synthesis. Unlike learning-based approaches, which may struggle with out-of-distribution terrains, our method remains adaptable to new conditions.

\textbf{Paragraph 1: We tackle terrain-adaptive locomotion. Problem 1: existing work lacks formal guarantee, but it is important for safety-critical applications.}
% Problem 2: existing MIP is computationally expensive.}

Terrain-adaptive locomotion is crucial for enhancing the traversing capabilities of legged robots and advancing beyond blind walking \cite{di2018dynamic,kim2019highly, zhou2022momentum}. Existing approaches often use heuristics \cite{fankhauser2018robust,jenelten2020perceptive,agrawal2022vision,grandia2023perceptive, asselmeier2024hierarchical} to simplify foot placement selection given a nominal reference trajectory. However, formal guarantees on locomotion safety, considering the robot's physical capabilities, are rarely explored—despite their importance in safety-critical scenarios like hazardous debris and construction sites. \textbf{Talk more about other approaches.}

\textbf{Paragraph 2: One popular way to provide formal guarantee is to formulate a big MIP. 
But it is computationally expensive, has scalability issue.
can have certain level guarantee.}
Moreover, complex features of surrounding terrains significantly increase the computational burden for optimization-based methods, such as mixed-integer program (MIP) \cite{valenzuela2016mixed,shirai2022simultaneous,jiang2023locomotion}, which directly address the interplay between terrain segments and robot kinematics and dynamics. Other MIP-based methods rely on simplifications to enhance speed, \textit{e.g.}, assuming predetermined contact sequences and timing \cite{deits2014footstep,ponton2016convex,ding2020kinodynamic,tonneau2020sl1m,risbourg2022real,acosta2023bipedal} or fixing the number of footsteps \cite{aceituno2017simultaneous}, but still face challenges in scaling up with the number of terrain segments.

\textbf{Paragraph 3: We propose an integrated framework that combines synthesis + MICP to solve (locomotion with formal guarantee), faster than MIP solution. Two-stage framework.}

Motivated by this fact, our study aims to combine formal control methods, \textit{i.e.}, linear temporal logic (LTL)-based reactive synthesis \cite{bloem2012synthesis}, with MIP-based approaches, to safely and promptly react to dynamically changing environments. Continuous robot states and local environments are abstracted into symbolic states to systematically study high-level interactions between the robot and surrounding challenging terrains such as obstacles and large gaps exceeding the robot's physical limit. Then, a Mixed-Integer Convex Program (MICP) is solved to provide a physical feasibility certificate for each symbolic transition. This integrated framework not only bridges the gap between high-level abstraction and low-level physical capabilities but also alleviates the computational burden of MICP through symbolic-level guidance, focusing on a nominal desired direction and specific regions to explore.
% focusing on a nominal desired moving direction and limiting the number of regions to explore. 
The proposed approach is flexible in that it can be used to follow commands from a global/higher-level planner \cite{wermelinger2016navigation, fernbach2017kinodynamic,norby2020fast,chignoli2022rapid,liao2023walking,feng2023gpf} or from a user, and can incorporate additional locomotion behavior as needed.

\textbf{Paragraph 4: Our offline stage has two key scalability issues. 1: synthesis: exponential time complexity wrt vars. Our approach includes terrains, and prohibitively large state space. 2: gait-free MICP is still expensive, so we want to minimize solving MICP.}

\textbf{Paragraph 5: 1. to solve expensive synthesis, we use a high-level manager to reduce the size of the state space, by fixing the values 
2. To solve the expensive MICP problem, we leverage symbolic repair to guide the skill generation, only generate necessary ones, reduce the number of times to run the expensive MICP}
Enabling complex reasoning about physical interaction with a terrain-feature-rich environment requires an accurate abstraction of the robot states or terrain. However, this induces a scalability challenge, particularly when evaluating numerous symbolic transitions for underactuated, hybrid systems such as legged robots. Rather than solving for abundant feasible transitions \cite{bhatia2010sampling,liu2013synthesis,maly2013iterative,gu2022reactive} for each terrain combination that potentially involve complex contact selection, it is more computationally efficient to utilize existing contact plans such as a fixed set of locomotion gait configurations, and ensure traversability among a limited set of terrains selected by prior knowledge or experience. In addition, merely deploying the existing proven feasible locomotion gaits provides another layer of safety.
% it is more computationally efficient to utilize existing motion primitives such as a fixed set of gait configurations, tailored to a limited set of terrains informed by prior knowledge or experience. In addition, deploying the existing proven feasible motion primitives provides another layer of safety. 
However, missing enough transition capabilities
% it 
often leads to \textit{unrealizable specifications}. Therefore, it is essential to develop a method that automatically identifies physically feasible symbolic transitions necessary for legged robots to maneuver in challenging environments. In this work, we adopt a symbolic repair approach based on \cite{pacheck2022physically}, incorporating dynamic feasibility checks, to design missing yet necessary locomotion gaits for our terrain-adaptive locomotion problem.

\textbf{Paragraph 6: Our online stage has two issues. We solve them by using MICP + repair online.}
During online execution of a synthesized high-level transition strategy, two critical issues may arise, potentially leading to execution failures. First, discrepancies between offline-checked terrain and real-world one are often inevitable, and can lead to a mismatch between the desired symbolic transition and what is physically possible. Second, newly encountered terrains may also make the symbolic specifications unrealizable, as they were not accounted for during offline synthesis. %Existing efforts in the LTL synthesis literature address unexpected environmental changes and robot action failures mainly at the specification level \cite{livingston2012backtracking,ehlers2014resilience,guo2013revising,li2021reactive,zhou2022reactive,cao2022leveraging,meng2024automated}. Studies that consider both robot and environmental mismatches at symbolic and physical levels remain largely underexplored, yet become increasingly crucial for safety-critical, terrain-adaptive locomotion. 
In this work, we solve an MICP online while executing the synthesized transition strategy, using real-world terrain data to ensure motion feasibility and fidelity. 
Moreover, 
we leverage online symbolic repair to
% an online symbolic repair is performed, similar to the offline repair, to 
identify possible transitions based on the new terrain and updated feasibility provided by the online MICP, allowing the robot to explore previously unseen terrains.
% for progressive generalization to a wider range of terrains.

\textbf{Contribution:
1. we combine synthesis + MICP to provide guarantee and reduce computational burden.\\
2. We use high-level manager to solve the scalability issue of synthesis,
use symbolic repair to reduce the computational burden of MICP.\\
3. Online repair }

The primary contributions of this paper are as follows:
\begin{itemize}
    \item We propose a hierarchically integrated planning framework for terrain-adaptive locomotion by combining reactive synthesis with MICP. This framework generates locomotion gaits, a computationally expensive process, based on the actual needs of the robot, by leveraging symbolic search over an abstraction to limit the number of required locomotion gaits.
    %\item To address unrealizable specifications due to limited predefined motion primitives, we introduce a symbolic repair approach that incorporates dynamic feasibility checks, and automatically identifies the missing but necessary transitions for legged robots to effectively navigate adversarial terrains.
    \item During online execution, we address the disparity between offline synthesis and real-world terrain conditions at both physical and symbolic levels. Our approach leverages an online MICP solver along with an online symbolic repair process to account for real-world terrain discrepancies and newly encountered conditions, enhancing motion feasibility and safety during execution.
    % \item The entire framework is implemented in an online manner and validated through both simulation and hardware experiments. To the best of our knowledge, this is the first real-time local planner that determines both gait and foothold with a safety guarantee.
\end{itemize}

% \section{Related Work}
\subsection{Hierarchical Planning for Terrain-Adaptive Locomotion}
Planning for terrain-adaptive locomotion can be boiled down into solving contact sequence and timing (equivalently gait), location (equivalently foothold), and robot motion (e.g., Center of Mass (CoM) or base pose), given the chosen kinematics/dynamics model and the perceived terrain information. Hierarchical approaches first separately or simultaneously address part of the above problems and then solve for the rest, which allows for higher computational efficiency.  
% Among them, hierarchical approaches commonly separate the foothold from the gait/contact and base pose planning, which allows for higher computational efficiency.
Kinematic footstep planning focuses on the first two problems, neglecting the robot dynamics through graph-based approaches \citep{kolter2008control,kalakrishnan2010fast,winkler2014path,mastalli2015line,fankhauser2018robust,griffin2019footstep} or optimization-based approaches \citep{deits2014footstep,tonneau2020sl1m,risbourg2022real, ren2025accelerating}.
% Another perspective of hierarchical approaches is on planning over a set of actions\citep{kolter2008control,kalakrishnan2010fast,winkler2014path,mastalli2015line,griffin2019footstep}.
Although non-fixed gait/contact plans can be generated for very rough terrains \citep{hauser2005non,bretl2006motion,tonneau2018efficient}, conservative quasi-static motions are generated for challenging terrains due to the kinematic simplification during  footstep planning. 
% While dynamically-feasible contact plans are also reported in, heavy computational time and difficulty in handling collision avoidance restrict the online re-planning.

Given the recent advances in optimization-based approaches, notably Model Predictive Control (MPC), another focus of the literature is on the local foothold adaptation and the integration of robot dynamics for dynamic locomotion.
Instantaneous foothold adaptation around a nominal location, with a fixed and cyclic gait, is performed based on heuristic search \citep{jenelten2020perceptive,kim2020vision,agrawal2022vision} or learning approaches \citep{magana2019fast, villarreal2020mpc,yu2021visual,gangapurwala2022rloc, wu2025learn, Kamohara2025RLMPC}. 
In \citep{grandia2021multi,jenelten2022tamols,grandia2023perceptive}, the base pose and foothold are optimized jointly, demonstrating impressive terrain traversing capabilities on rough terrains. More recently, MPC has been employed within hierarchically integrated planning frameworks for legged navigation over rough terrain, incorporating explicit terrain uncertainty quantification \citep{muenprasitivej2024bipedal, shamsah2025terrain, muenprasitivej2026bipedal}.
% , which generates dynamically consistent trajectories. 

However, the above dynamic locomotion works consider a fixed gait pattern that usually prohibits achieving versatile behaviors, such as jumping.
In \citep{Park2015OnlinePF, kim2020vision,gilroy2021autonomous}, jumping motions are generated offline and triggered through heuristics.
In \citep{fernbach2017kinodynamic,norby2020fast,chignoli2022rapid}, RRT-Connect is used for rapid kino-dynamic locomotion planning, and the switch between normal walking and jumping is either embedded inside a reduced-order model during sampling \citep{norby2020fast} or decided by a pre-trained feasibility classifier \citep{chignoli2022rapid}.
Beyond limited gait modes, other works focus on online gait planning and/or foothold selection and deploy MPC to track the plan,
which can be further categorized into model-free \citep{da2021learning,yang2022fast,xie2022glide,pmlr-v164-margolis22a,yu2024learning} and model-based methods \citep{boussema2019online,wang2023and,sun2023free}.

\revised{A complementary line of work uses end-to-end reinforcement learning (RL) to train perceptive locomotion policies that directly map terrain observations to joint commands, achieving impressive robustness on stairs, gaps, and unstructured outdoor terrain~\citep{miki2022learning,agarwal2023legged,zhuang2023robot,hoeller2023anymal,cheng2024extreme,luo2024pie}. Compared with these data-driven approaches, our model-based, formal-methods framework offers explicit symbolic-level realizability guarantees and physical-feasibility certificates from MICP, at the cost of weaker generalization and the need for a terrain abstraction; we view the RL approaches and ours as complementary rather than competing.}

\subsection{Simultaneous Planning via Optimization}
Planning for terrain-adaptive locomotion problem can also be formulated as a combinatorial optimization problem that simultaneously optimizes for gait, foothold, and robot dynamics.
One approach is to incorporate rigid \citep{posa2014direct} or smoothed \citep{mordatch2012discovery, drnach2021robust} complementarity constraints, which accurately model the physical contact interaction behavior but remains computationally inefficient due to the non-smoothness and high nonlinearity. Promising online contact-implicit MPC results are reported in \citep{aydinoglu2022real,le2024fast, drnach2022mediating} but still limited to low-dimensional system or static environment. 

Another way of encoding discrete contact decisions is to treat both the contact state and the contact plane selection for each leg at each timestep as binary variables \citep{valenzuela2016mixed,shirai2022simultaneous,jiang2023locomotion}. The underlying problem can be formulated as Mixed Integer Program (MIP). Convex approximations of nonlinear dynamics, usually centroidal dynamics, and constraints are typically required to transform the MIP into a more tractable Mixed Integer Convex Program (MICP), albeit with an increased number of binary variables. 
A global certificate for the approximated convex problem exists upon convergence \citep{dai2019global}, providing an ideal means to determine the feasibility in our proposed method.
To relieve the computational burden due to the introduction of many binary variables, the works of \citep{ponton2016convex,ding2020kinodynamic,acosta2023bipedal} assume the contact sequence and timing are chosen \textit{a priori}, and only optimize the contact plane selection. Aceituno-Cabezas et al. \citep{aceituno2017simultaneous} optimize the contact sequence within a certain gait cycle that fixes the number of footsteps. 
However, it is inevitable that the problem complexity grows exponentially when the number of discrete contact options increases along with the time horizon and terrain features. Our work aims to alleviate computational burden through a two-stage approach. In the offline phase, expensive MICPs that optimize both contact state and foothold selection are solved to generate necessary locomotion gaits. In the online phase, efficient MICPs with adaptive and predefined gaits are used to produce dynamically feasible motions and footholds. Additionally, guided by a synthesis-based task planner, each MICP in our work considers shorter time horizons and fewer terrain features compared to solving a single large MICP problem.

% Our approach: between hierarchical planning and combinatorial optimization. Split the contact planning problem into gait and footstep selections, while the gait planner considers dynamical feasibility given the terrain information.

% Different from \citep{shirai2022simultaneous}, this work: 

% 1. Considers the angular evolvement of the robot to allow more versatile motion generation 

% 2. We consider a more complicated centroidal optimization problem which provides higher-fidelity solution back to the MICP 

% 3. We fix the gait selection to allow real-time foothold selection 

 

% Different from the LGP work \citep{toussaint2015logic}: 

% 1. We evaluate the entire course of a trajectory and incorporate the underactuated dynamics, which is hard to be divided into multiple sub-problems. 

% Different from the MCTS \citep{amatucci2022monte}: 

% 1. We mostly care about the computational efficiency and the motion feasibility across a certain course. Therefore, two different choices we made here: 

% 2. Our footstep planner assumes a fixed gait and only determines the foothold locations given the steppable regions. 

% 3. The gait selection in our problem setup stays in a relatively higher level, whereas the MCTS is able to discover and optimize the gait. In other words, the high-level search space in our framework is smaller than MCTS.   

% Independent Gait Scheduling

% As discussed above, simultaneously planning for the gait and the footholds given perception feedback remains a challenging problem. Alternatively, impressive results have been reported by separating the gait scheduling and the foothold selection. The state-of-the-art works that handle gait transitions can be divided into model-free \citep{da2021learning,yang2022fast,xie2022glide,pmlr-v164-margolis22a}... and model-based approaches \citep{boussema2019online,wang2023and,sun2023free}.

\subsection{Task and Motion Planning for Contact-Rich Planning}
% Instead of capturing every discrete option as in general MIP, 
Task and Motion Planning (TAMP) formally defines symbolic-level tasks and searches through a graph of predefined motion primitives that enable feasible symbolic transitions \citep{garrett2021integrated, zhao2024survey}. For more complex physical contact resoning, Toussaint et al. propose Logic Geometric Programming (LGP) \citep{toussaint2015logic} and embed the high-level
logic representation into the low-level motion planner, demonstrating contact-rich tool-use behaviors \citep{toussaint2018differentiable}
after defining abundant action primitives. 
Building on this, a broader range of motions has been showcased, including versatile manipulation \citep{migimatsu2020object,zhao2021sydebo} and loco-manipulation tasks \citep{sleiman2023versatile}.
In a similar vein, works such as \citep{ding2021hybrid,shirai2021lto,jelavic2021combined,amatucci2022monte,chen2021trajectotree,zhang2023simultaneous,zhu2023efficient,asselmeier2024hierarchical} integrate graph-search or sampling-based methods with optimization-based approaches in a holistic manner, abstracting contact modes within a discrete domain. However, the combinatorial nature of these problems often leads to poor scalability due to the explosion of contact modes. Deploying such approaches online in dynamic environments remains an open challenge. From a different perspective, we abstract a smaller set of task-level contact modes as locomotion gaits over a specific time horizon and distance, allowing the MICP to solve for details such as contact locations and robot motion during execution. This shifts the focus to selecting locomotion gaits within a local environment while ensuring safety and completeness through synthesis-based approaches.

Unlike LGP that solves an integrated TAMP problem, synthesis-based approaches using Linear Temporal Logic (LTL) operate primarily at the task level, emphasizing safety and completeness guarantees within the task domain \citep{kress2018synthesis}. These methods typically synthesize a sequence of task-level actions, which are then executed by a motion planner. Recently, LTL has been applied to safe locomotion tasks, employing distinct task-level abstractions on topological maps of terrain \citep{kulgod2020temporal,zhou2022reactive,jiang2023abstraction} or locomotion keyframes for reduced-order models \citep{gu2022reactive,shamsah2023integrated,zhao2022reactive}.
Reactive synthesis \citep{pnueli1989synthesis}, widely applied to mobile robots \citep{kress2009temporal}, has been employed to enable prompt decision-making in response to more complex environments \citep{shamsah2023integrated} or external perturbations \citep{gu2022reactive}. Notably, \citep{zhao2022reactive} formulate the terrain-adaptive locomotion problem as a sequential decision-making task, selecting appropriate locomotion modes with predefined contact and locomotion keyframes to accommodate dynamically changing terrains. However, the locomotion mode selection is explicitly encoded in the LTL specification using expert knowledge, lacking a physically feasibility guarantee.
% To reduce the complexity of solving multi-contact, long-horizon, and kino-dynamic planning problems, a single optimization problem can be divided into multiple sub-problems, facilitating faster computation through holistic approaches. Graph-search planning (GSP) or sampling-based planning (SBP) methods have been proposed to be combined with dynamically feasible motion primitives \citep{shkolnik2011bounding, fernbach2017kinodynamic,norby2020fast,shirai2021lto,ding2021hybrid,jelavic2021combined,chignoli2022rapid,amatucci2022monte}. Superior planning speed has been reported with reduced-order models \citep{fernbach2017kinodynamic,norby2020fast,carpentier2018multicontact} or pre-trained feasibility classifier \citep{chignoli2022rapid}. 
% The authors in \citep{zhao2021sydebo,shirai2021lto,ding2021hybrid,jelavic2021combined,amatucci2022monte,chen2021trajectotree,zhang2023simultaneous} fully exploit the robot dynamics through TO combined with a high-level search. 
% Among them, researchers in \citep{jelavic2021combined,amatucci2022monte} explicitly reason the gait/contact mode in the search process. 
% However, the combinatorial nature of the problem renders poor scalability due to the explosion of contact modes. Online deployment of this type of approaches given changing terrain information remains an open problem.
In this work, we aim to combine the strengths of reactive synthesis and MICP, utilizing the global certificate of MICP as a feasibility checker for terrain-adaptive locomotion.

% The lack of immediate reaction and safety guarantee given real-time terrain information still prohibits this type of approach from online deployment.

\subsection{Physically-Feasible Reactive Synthesis}
% Linear Temporal Logic (LTL) serves as a prevalent tool for encoding temporal task specifications and automating the synthesis of system transition behaviors within the robotics community. Originating from its proposal for model checking by Pnueli in 1977 \citep{pnueli1977temporal}, LTL has gained substantial traction in robotics research, initially prevalent in studies involving wheeled robots \citep{kress2009temporal, decastro2018collision}, and more recently, extending to legged robots \citep{warnke2020towards, zhao2022reactive, kulgod2020temporal, gu2022reactive, shamsah2023integrated}.
% Reactive Synthesis, denoting the automated generation of a reactive system to ensure compliance with desired specifications, regardless of external inputs, was introduced by Pnueli in 1989 \citep{pnueli1989synthesis}. Given the escalating complexity associated with lengthy LTL formulas, a subset of LTL, General Reactivity of Rank 1 (GR(1)), was proposed to streamline the synthesis process \citep{piterman2006synthesis}. This subset proves especially beneficial in managing the synthesis process efficiently.

% (Replanning in the task level; BT, TAMP; inflation? Check Jonas paper again) 

% In contrast to Task and Motion Planning (TAMP), a paradigm that typically integrates symbolic search with geometric or optimization-based motion planners, reactive synthesis primarily operates at the task level. 

% Reactive synthesis primarily operates at the task level and involves specifying Linear Temporal Logic (LTL) formulas that encompass sequencing, conditions, and reactions to environmental events. 
While reactive synthesis can promptly and safely respond to dynamic environments, it typically does not assess physical feasibility when being deployed on complex robotic systems. To bridge the gap between discrete abstraction and continuous system dynamics, various strategies have been proposed. Studies in  \citep{decastro2014dynamics,fainekos2006translating} focus on automatically synthesizing controllers in the continuous domain based on high-level specifications. Another approach involves incorporating dynamics directly into the reactive synthesis process by abstracting physical systems, including nonlinear \citep{bhatia2010sampling}, switched \citep{liu2013synthesis}, and hybrid systems \citep{maly2013iterative} with manageable model complexity. However, for legged robots with multiple contacts navigating dynamic terrains, these physically feasible reactive synthesis approaches remain computationally intractable.

Recent efforts \citep{pacheck2020finding, pacheck2022physically} focus on symbolic repair for unrealizable task specifications, defining symbolic skills with preconditions and postconditions, similar to TAMP, to identify missing skills that make the specification realizable. These works introduce iterative feasibility checks based on symbolic repair suggestions, emphasizing alignment between symbolic task specifications and physical reality. Building upon this, Meng et al. \citep{meng2024automated} extend this type of approach to online symbolic repair.
Inspired by these works, Zhou et al. \citep{zhou2025physically} adopt a similar mechanism for terrain-adaptive locomotion. Specifically, they abstract the robot and terrain states in a local environment and define physically feasible skills by solving MICP. 
% To address scalability issues arising from complex terrain features, they employ a high-level manager to decompose the problem into smaller subproblems, enabling faster repair and synthesis. 
They leverage high-level manager and symbolic repair to efficiently identify missing but necessary skills, reducing the need for frequent calls to the computationally expensive gait-free MICP. Our work significantly completes \citep{zhou2025physically} through seamless integration with a low-level tracking control module, enabling systematic benchmarking, and demonstrating real-world hardware deployment.

% Three main novelties in our reactive planner:

% 1) We leverage the powerfulness of reactive synthesis to  reason about the robot traversability with regard to the terrain information.

% 2) We generate skills that are physically feasible offline; a symbolic repair process is proposed to handle unrealizable specifications.

% 3) We perform fast run-time feasibility checking through graceful coordination toward MPC given real-time robot and environment representations.

\section{Preliminaries}
Let us use an example to elaborate our framework:
% throughout the paper.

\noindent
\textbf{Example~\examplelabel{example:1}}.
% (9-grid locomotion). 
\textit{Consider the 2D grid world with nine cells in Fig.~\ref{fig:preliminary_nine_squares}. The robot must move from an initial square (e.g. the center square) 
to a desired goal square indicated by the star. 
Each cell is assigned a terrain type, such as $\textit{Dense}$ or $\textit{Sparse}$ stepping stones.}

%%%%% Fig. 9 grid %%%%
\begin{figure}[t]
    \centering
    \includegraphics[width=\linewidth]{preliminary_drawing.pdf}
    \caption{Terrain abstraction and skill definition. (a) Top-down view of the real-world terrain before abstraction. The red polygons denote the segmented terrain polygons. (b) Abstraction of the terrain and robot's skills (pink) moving from one location to another. (c) Repair process to find a new skill (blue).}
    \label{fig:preliminary_nine_squares}
\end{figure}
%%%%%%%%%%%%%%%

\subsection{Abstractions}\label{sec:preliminaries_abstraction}
We abstract the environment states and the robot actions into a set of atomic propositions $\ap$.
We partition $\ap$
% $ = \mathcal{I} \cup \mathcal{O}$ 
into sets of inputs $\mathcal{I}$
% $= \{\mathcal{X}_{\rm{cen}}, \mathcal{E}, \mathcal{X}_{\rm{req}}\}$ 
% (``environment" in the two-player game) 
and outputs $\mathcal{O}$
($\ap = \inp \cup \out$).
% (``robot" in the two-player game). 
The inputs $\mathcal{I}$ consist of
the centroidal inputs $\inpcen$,
% which describe the robot's centroidal states,
% $\mathcal{X}_{\rm{cen}}$ of the robot, 
the terrain inputs $\inpterrain$,
% that describe the environment states surrounding the robot,
% $\mathcal{E}$ that surrounds the robot, 
and the request inputs $\inprequest$
% that specify the goal states provided by the higher-level planner or the user
% the requested robot state $\mathcal{X}_{\rm{req}}$ given by a higher-level planner or user input
($\inp = \inpcen \cup \inpterrain \cup \inprequest$).
The centroidal inputs $\inpcen$ describe the robot's centroidal states,
grounding to the robot's centroidal keyframes
defined as $\mathcal{X} = \{x, y\}$, where the translational states $x$ and $y$ are located at the center of each grid based on the Cartesian-space discretization. 
% A centroidal state $\inpstatecen \subseteq \inpcen$
% We exclude the vertical height of the robot since it can be purely decided by the terrain information. 
%
The terrain inputs $\inpterrain$
characterize the terrains around the robot,
providing abstract information about the terrain,
e.g. their heights
or densities.
% terrain sparsities, or obstacle detection.
% given the classified terrain information: $\mathcal{E} = \{\rm{cT}, \rm{lT}, \rm{rT}, \rm{uT}, \rm{dT}\}$, where $\rm cT$ stands for current terrain, $\rm lT$ stands for left terrain, $\rm rT$ stands for right terrain, $\rm uT$ stands for upper terrain, and $\rm dT$ stands for down terrain. 
% Each environmental input encodes abstracted information about the terrain, 
% such as the height, sparsity, 
% or obstacle detection. 
%
The request inputs $\inprequest$ allow the high-level planner or the user to specify the target robot states.
%
Since the robot can only control the values of centroidal inputs $\inpcen$, 
we also refer to centroidal inputs $\inpcen$ as controllable inputs
and
terrain and request inputs ($\inpterrain\cup \inprequest$) as uncontrollable inputs.
%
In Example~\ref{example:1},
$\inpcen \coloneqq \{x_0, x_1, x_2, y_0, y_1, y_2\}$,
$\inprequest \coloneqq \{v^\textrm{req} \mid v \in \inpcen\}$,
and
$\inpterrain \coloneqq \{\tvar{x_i}{y_j} \mid i,j\in \{0, 1, 2\}\}$
where each input $\tvar{x_i}{y_j}$ is True iff the terrain $(x_i, y_j)$ is Dense.
%
% A centroidal state $\inpstatecen \subseteq \inpcen$
% is a 
We define 
the
centroidal states $\inpstatecen \subseteq \inpcen$, 
terrain states $\inpstateterrain \subseteq \inpterrain$, 
and request states
$\inpstaterequest \subseteq \inprequest$
as
subsets of 
centroidal inputs $\inpcen$, 
terrain inputs $\inpterrain$, 
and request inputs $\inprequest$,
respectively.
% Thus is the same as the robot states ($\inprequest = \inpcen$).
%
% The requested state $\mathcal{X}_{\rm{req}}$ is defined same as the robot state: $\mathcal{X}_{\rm{req}} = \mathcal{X}_{\rm{cen}}^{\rm{des}}$. 
% One example with only translational centroidal states is shown in Fig.~\ref{fig:cartoon_problem}(a).
%
In Fig.~\ref{fig:preliminary_nine_squares}b,
the centroidal state is $\{x_1, y_1\}$,
the request state is $\{x_0^{\textrm{req}}, y_2^{\textrm{req}}\}$,
and the terrain state is $\{\tvar{x_1}{y_0}, \tvar{x_1}{y_1}, \tvar{x_1}{y_2}\}$.

% \qm{To be deleted} Furthermore, we partition the inputs $\inp$ into controllable inputs $\inpsyms$ and uncontrollable inputs $\inpuser$.
% Controllable inputs $\inpsyms$ represent the environment states the robot can control, i.e. the robot's centroidal states ($\inpsyms = \inpcen$).
% Uncontrollable inputs describe the part of the physical world that the robot does not have control over, i.e. 
% the terrain states $\inpterrain$ and the requested goal states $\inprequest$ ($\inpuser = \inpterrain \cup \inprequest$).

The outputs
% aforementioned output propositions 
$\out$ 
represent
% consists of 
the skills 
that allow the robot to transit among centroidal states,
given terrain states.
% $\mathcal{A}$, which define the capabilities of transitioning from one keyframe $\mathcal{X}_{\rm cen}$ to another given certain environment $\mathcal{E}$. 
% Following a TAMP-style definition for future repair process, 
A skill $\skill \in \out$
consists of 
a set of preconditions $\skillpre \subseteq 2^{\inpcen \cup \inpterrain}$
from which the skill is allowed to execute,
and a set of postconditions $\skillpost \subseteq 2^\inpcen$, the resulting centroidal states after executing the skill.
%
In Example~\ref{example:1},
the skill $\skillcustom{0}$ moves the robot from the middle to the upper terrain grid (Fig.~\ref{fig:preliminary_nine_squares}b).
The precondition of $\skillcustom{0}$ is
$\skillprecustom{\skillcustom{0}} = \{\{x_1, y_1, \tvar{x_1}{y_0}, \tvar{x_1}{y_1}, \tvar{x_1}{y_2}\}\}$, and
the postcondition is
$\skillpostcustom{\skillcustom{0}} = \{\{x_0, y_1\}\}$.
% More details of this skill will be provided in Sec.~\ref{sec:spec_gen}.
Each skill not only contains the abstraction-level information for defining the transition, but also corresponds to a continuous-level locomotion gait, such as a one-second trotting gait, which is further defined in Sec.~\ref{sec:motion_primitive}. 

%%%% Fig.2 System architecture %%%%
\begin{figure}[t]
    \centering
    \includegraphics[width=1\linewidth]{system_diagram.pdf}
    \caption{System architecture overview. The solid lines indicate online communication, while dashed lines represent offline processes.
    }
    \label{fig:system_diagram}
\end{figure}
%%%%%%%%%%%%%%%%%%%%%%%%%%%%%%%%%%%

% \subsection{Reactive  Synthesis}
\subsection{Specifications}
We use the Generalized Reactivity(1) (GR(1)) fragment of Linear Temporal Logic (LTL)~\cite{bloem2012synthesis}
% has been widely used 
to encode task specifications.
% ~\cite{bloem2012synthesis}.
% and automatically synthesize the system's transition behaviors. 
GR(1) specifications $\spec$ 
% is constructed from atomic propositions $\pi \in AP$, which is evaluated to be ${\rm True}$ or ${\rm False}$ and 
follow the grammar 
$\spec \: := \var \:|\: \neg\: \spec \:|\: \spec \:\land\: \spec \:|\: \bigcirc\spec \:|\: \square \spec \:|\: \lozenge \spec$,
where $\var \in \ap$ is an atomic proposition,
$\neg$ ``not'' and $\land$ ``and'' are Boolean operators,
and
$\bigcirc$ ``next'',
$\square$ ``always'',
and $\lozenge$ ``eventually''
are temporal operators.
% . Boolean operators $\neg$ ``not'' and $\wedge$ ``and'' in addition to a set of temporal operators $\ocircle$ ``next'', $\mathcal{U}$ ``until'', and $\mathcal{R}$ ``release'', are denoted. To be concise, we omit the derivations of other boolean operators such as $\vee$ ``or'', $\rightarrow$ ``implies'', and $\leftrightarrow$ ``if and only if'', as well as temporal operators $\Diamond\,\varphi$ ``eventually $\varphi$'' and $\square\,\varphi$ ``always $\varphi$''. 
We refer the readers to~\cite{baier2008principles} for a detailed description of LTL.
% A detailed description of LTL can be referred to \cite{baier2008principles}.

GR(1) specifications are expressed in the form of 
$\spec = \envspec \to \sysspec$,
where $\envspec = \envinit \land \envsafety \land \envlive$ 
is the assumptions on the behaviors of the possibly adversarial environment,
and $\sysspec = \sysinit \land \syssafety \land \syslive$
is the guarantee of the robot's behaviors.
% Specifically,
For $\alpha \in \{e, s\}$,
$\alphainit, \alphasafety, \alphalive$
characterize the 
initial conditions,
safety constraints,
and liveness requirements
respectively.
%
We divide the environment safety assumptions $\envsafety$ into skill assumptions $\envsafetynothard$,
which describe the postconditions of the skills (see Sec.~\ref{sec:spec_gen}),
and hard assumptions $\envsafetyhard$ that cannot be modified
($\envsafety = \envsafetynothard \land \envsafetyhard$).
Similarly, we split the safety guarantees $\syssafety$ into skill guarantees $\syssafetynothard$,
which represent the preconditions of the skills,
and hard guarantees $\syssafetyhard$ that cannot be modified
($\syssafety = \syssafetynothard \land \syssafetyhard$).



% In this work, we consider specifications expressed in the Generalized Reactivity(1) (GR(1)) fragment of LTL \cite{bloem2012synthesis}. The problem of realizability is explored in a two-player game between the robot and its possibly adversarial environment. The winning condition, as known as the GR(1) condition, is provided by:
% \begin{equation}
    % \varphi = \varphi_e \rightarrow \varphi_s    
% \end{equation}
% where $\varphi_e = \varphi_e^i \wedge \varphi_e^t \wedge \varphi_e^g$ specifies the environment behavior and $\varphi_s = \varphi_s^i \wedge \varphi_s^t \wedge \varphi_s^g$ specifies the system's behavior. More specifically, $\varphi^i,\varphi^t$ and $\varphi^g$ characterize the initial states, safety constraints, and liveness requirements, respectively.

% \subsection{9-Grid Example for Terrain-Adaptive Locomotion}
% For the ease of elaborating our algorithm, let us introduce a basic 9-grid locomotion example first. 

% \begin{figure}[t]
%     \centering
%     \includegraphics[width=\linewidth]{preliminary_drawing.pdf}
%     \caption{Terrain abstraction and skill definition. (a) Top-down view of the real-world terrain before abstraction. (b) Abstraction of the terrain and robot's skills (pink) moving from one location to another. (c) Repair process to find a new skill (blue).}
%     \label{fig:preliminary_nine_squares}
% \end{figure}

% % In this work, we consider a specific perceptive locomotion problem, where
% \noindent\textbf{Example~\examplelabel{example:1}}
% (9-grid locomotion). 
% \textit{Consider a 2D grid world with 9 discretized grid cells in Fig.~\ref{fig:preliminary_nine_squares}. The robot is commanded to move from an initial square (usually the center square), to a desired goal square indicated by the star. Each grid is assigned with a terrain type, such as $\textit{Dense}$ and $\textit{Sparse}$ stepping stones.}

% %The terrains are only characterized into $\textit{Flat}$ and $\textit{High}$ ones, although 
% More complex terrain types can also be defined. In Sec.~\ref{sec:experiments}, we will show more examples with realistic set-ups.

\section{Problem Statement}\label{sec:problem_statement}
% Given a set of inputs $\inp = \inpcen \cup \inpterrain \cup \inprequest$ after discretizing the local environment and categorizing the terrain polygons $\mathcal{P}$, and a set of locomotion primitives, 
We aim to solve a decision-making problem of terrain-adaptive locomotion.

% that generates a symbolic strategy $\mathcal{A}$, consisting of physically feasible skills, to safely and reliably navigate the local environment. To achieve this, the following problems must be addressed, spanning symbolic abstraction to motion planning:

\textbf{Problem 1}.
% (Sec.~\ref{sec:preliminaries_abstraction} and Sec.~\ref{sec:reactive_gait_synthesis}):} 
Given a set of predefined request goals and well-structured terrain polygons, 
find a symbolic strategy $\strategy$, consisting of physically feasible skills, to safely and reliably navigate the local environment
% abstract them into a set of inputs $\inp = \inpcen \cup \inpterrain \cup \inprequest$, find necessary and physically-feasible skills and corresponding locomotion gaits, encode and synthesize realizable task specifications.

\textbf{Problem 2}.
% (Sec.~\ref{sec:online_execution}):} 
Given the synthesized strategy $\mathcal{A}$, an online request goal, and online segmented terrain polygons, 
address 
% potential online execution failures, including 
real-world terrain discrepancies and unseen terrain types during execution.

% \section{System Architecture}
\section{Approach Summary}
To tackle the above problems, our planning framework consists of two main modules: 
offline synthesis (Sec.~\ref{sec:reactive_gait_synthesis}) 
and
online execution (Sec.~\ref{sec:online_execution}), as shown in Fig.~\ref{fig:system_diagram}.
% a MICP planner (Sec.~\ref{sec:motion_primitive}),
% and a tracking controller (Sec.~\ref{sec:tracking_control}). 

The offline synthesis module generates a strategy that moves the robot to a set of predefined goal waypoints among possible terrain polygons,
utilizing a MICP planner to obtain feasible robot skills corresponding to the symbolic transition.
% If the skills are not enough to generate a strategy for the robot to reach any goal waypoint over any possible terrain states, we leverage symbolic repair to generate new skills and utilize the gait-free MICP planner to check the feasibility of the suggestions.
Instead of considering predefined goals and terrains as done during offline, the online execution module
takes in an online local waypoint
produced by a global navigation planner,
along with the online segmented terrain polygons \cite{miki2022elevation}.
% a global navigation planner provides a series of local goal waypoints,
% and a perception module 
%
% supplies the segmented terrain information.
It then executes the strategy to navigate the robot toward the goal waypoint, 
utilizing the same MICP planner but with online terrain segments,
and leverages online repair 
% (Sec.~\ref{sec:online_repair}) 
to handle unexpected terrain configurations.

\section{Offline Synthesis}\label{sec:reactive_gait_synthesis}
The offline synthesis module 
% $\inpstateterrainposs \subseteq 2^\inpterrain$,
% and 
generates a strategy 
that enables the robot
to reach desired goal waypoints over all predefined terrain states.
%
As shown in Fig.~\ref{fig:combined_offline_online}(a),
this module takes in
a set of request goals and
a set of possible terrain states from prior knowledge of the workspace. In addition, a set of locomotion gaits is also provided by the user. After discretizing the local environment and characterizing the terrains in Sec.~\ref{sec:preliminaries_abstraction}, we first solve a gait-fixed MICP problem to determine the feasibility of each potential skill given the locomotion gaits
(Sec.~\ref{sec:motion_primitive}), and then encode all feasible transitions as skills in a specification
(Sec.~\ref{sec:spec_gen}).
If the specification is realizable, we generate a strategy for online execution.
Otherwise, we leverage a repair process to suggest new robot skills, utilize the gait-free MICP planner to check the feasibility of the suggested skills, and eventually make the specification realizable
(Sec.~\ref{sec:symbolic_repair}).


% The reactive synthesis process has two phases: offline synthesis and online execution. 
% We begin by introducing the essential specifications that define our terrain-adaptive locomotion problem. The objective is to construct a high-level robot gait and keyframe transition strategy by solving a two-player game between the gait planner and its possibly adversarial environment. The Linear Temporal Logic (LTL) specifications are designed by human experts, who provide logical rules that guide the planning process. It is nothworthy that the specification design is not unique and may vary according to expert input.

% Add a mapping function from the discrete state to the continuous state?

% \subsection{Motion Primitives to Robot Skills}\label{sec:motion_primitive}
\subsection{Locomotion Gait and Feasibility Checking via MICP}\label{sec:motion_primitive} 
We assume each locomotion gait comes with a contact sequence $\mathcal{G}$ and corresponding time durations $T$. Therefore, we define the locomotion gait as $L = \mathcal{M}(\mathcal{G}, T)$. Each skill is assumed to have a unique locomotion gait to determine how the robot moves at the continuous level. 
As shown in Fig.~\ref{fig:combined_offline_online}(a), since all the possible request goals $\inprequest$ and terrain states $\inpterrain$ are predefined during the offline phase, it is straightforward to first identify all the potential skills at the symbolic level that allow the robots to move freely in all possible local environments. 
% In most examples shown in this paper, we assume the robot is only allowed to move horizontally and vertically by one cell and we show diagonal movements and heading angle change in Sec.~\ref{sec:extension}.

\begin{figure}[t]
    \centering
    \includegraphics[width=\linewidth]{mip_go2.pdf}
    \caption{Demonstration of decision variables and polygons when solving MICP for skill $a_0$ in Example~\ref{example:1} using a trotting locomotion gait.}
    \label{fig:mip_go2}
\end{figure}

%%%% Fig.3 MICP formulation %%%%
\begin{figure}[ht]
    \centering
    
    % \boxed{
    \begin{minipage}{0.9\linewidth}
    \begin{subequations}\label{ocp:mip_general}
    \begin{alignat}{2}
        \nonumber &\underset{\vg \phi, \v H_{r,k,j}}{\min} \hspace{0.1cm} &&\sum_{i = 0}^{N-1} \hspace{2pt} \delta \vg \phi[i]^T\mathbf{Q}\hspace{2pt}\delta \vg \phi[i] + \vg \phi[i]^{T}\mathbf{R}\vg \phi[i]\\
     \nonumber & \text{ subject to} && \forall i \in \range{N}, j\in\range{n_f}, k\in\range{n_s}, r\in\range{R},\\
        \nonumber
        \\\label{eq:dynamics}
        &(\textit{\rm Dynamics}) &&  
        \begin{bmatrix} \v r[i+1]\\\
        \dot{\v r}[i+1]\\\
        m \ddot{\v r}[i]\\\
        \vg \theta[i+1]\\
        \dot{\vg \theta}[i+1]
        \end{bmatrix}= \begin{bmatrix}
        \v r[i] + \Delta t \cdot \dot{\v r}[i+1]\\
        \dot {\v r}[i] + \Delta t \cdot \ddot{\v r}[i+1]\\
        \sum\nolimits_{j} \v f_j[i] + m \v g \\
	    % \sum\nolimits_j (\v{p}_j[i] - \v r[i]) \times \v f_j[i] \\
            \vg \theta[i] + \Delta t \cdot \dot{\vg \theta}[i+1]\\
            \dot{\vg \theta}[i] + \Delta t \cdot \ddot{\vg \theta}[i+1]
	    \end{bmatrix}
        \\\label{eq:safe_region}
        &(\textit{\rm Foothold})
        && \v H_{r,k,j} \Rightarrow \v A_r \v p_j \leq \v b_r
        \\\nonumber
        & &&\sum_{r=0}^{R-1}\v H_{r,k,j} = 1
        \\\nonumber
        & &&\v H_{r,k,j} \in \{0, 1\}
        \\\label{eq:force_non_zero}
        &(\textit{\rm Frictional}) 
        % \hspace{1.05cm} \text{if} \ i \in \mathcal{C}_j \ \text{(in contact):}\\
        &&\v f_j[i] \cdot \v n(\v p_j^{xy}[i]) \ge 0, \ \forall i \in \mathcal{C}_j
        \\\label{eq:force_friction_cone}
        & &&\v f_j[i] \in \mathcal{F}(\mu, \v n, \v p_j^{xy}[i]), \ \forall i \in \mathcal{C}_j
         \\\label{eq:ee_velocity_zero}
            &(\textit{\rm Contact}) 
            &&\dot{\v p}_j[i]=0, \ \forall i \in \mathcal{C}_j
        \\\label{eq:force_zero}
         & && \v f_j[i] = 0, \ \forall i \notin \mathcal{C}_j
         \\\label{eq:max_torque}
         &(\textit{\rm Actuation})
         && \v J_j^{\top} \v f_j[i] \leq \vg \tau_{\rm max}
        \\\label{eq:rom}
            &(\textit{\rm Kinematics}) &&\v p_j[i] \in \mathcal{R}_j(\v r[i],\vg \theta[i])
    \end{alignat}
\end{subequations}
\end{minipage}
% }
    \caption{MICP formulation for physical feasibility checking.}
    \label{fig:micp}
    \vspace{-0.15in}
\end{figure}
%%%%%%%%%%%%%%%%%%%%%%%%%%%%%%%%%%%

% Eq.~\ref{ocp:mip_general} 
Fig.~\ref{fig:micp}
shows our MICP formulation, where the decision variables $\vg \phi$ include base position $\v r$, velocity $\dot{\v r}$, acceleration $\ddot{\v r}$, orientation $\vg \theta$ parameterized by Euler angle and its derivative $\dot{\vg \theta}$, individual 
end-effector
(EE)
position $\v p_j$, velocity $\dot{\v p}_j$, and acceleration $\ddot{\v p}_j$, and individual contact force $\v f_j$ for the foot $j$. The cost function consists of tracking costs and regularization terms governed by diagonal matrices $\v Q$ and $\v R$. 
The tracking costs include the deviation from a desired base trajectory interpolating a path between the terrain grid locations defined in each skill.
For any natural number $n \in \naturals$,
we denote the set $\{0, \dots, n-1\}$ as $\range{n}$.
We use $N$, $n_f$, and $n_s$, and $R$ to represent the number of timesteps, the number of feet, the number of footsteps specified by the gait configuration for each foot, and the number of convex terrain polygons to be considered, respectively. For example,  Fig.~\ref{fig:combined_offline_online}(b) shows a case with four polygons. Note that, during the offline phase, the terrain polygons are predefined and homogeneous, as depicted in Fig.~\ref{fig:preliminary_nine_squares}(a).
% \qm{plz also explain $N, n_f, n_s,$ and $R$}

% In Eq.~\ref{eq:dynamics},
We encode the system dynamics as a single rigid body in Eq.~\ref{eq:dynamics},
using a simplified Euler equation to keep the dynamics constraint convex. To incorporate safe region constraints to select proper footholds, we introduce binary variables $\v H_{r,k,j}$, with $r$, $k$, and $j$ expressing the $r^{\rm th}$ convex region, $k^{\rm th}$ footstep, and $j^{\rm th}$ foot. 
Eq.~\ref{eq:safe_region} restricts
the robot's EE to stay within one of the convex polygons. 
Frictional constraints are defined in Eqs.~\ref{eq:force_non_zero} - \ref{eq:force_friction_cone}, with $\v n$ as the normal vector of the terrain corresponding to the position $\v p^{xy}_j$ and $\mathcal{C}_j$ as the set of timesteps indicating stance for the $j^{\rm th}$ foot. Contact constraints in Eqs.~\ref{eq:ee_velocity_zero} - \ref{eq:force_zero} check the given contact states and activate at different timesteps.
Since the joint angles are omitted in the single rigid body model, we use a fixed Jacobian at a nominal joint pose to approximately consider the torque limit constraint in Eq.~\ref{eq:max_torque}. 
Lastly, the kinematics constraint in Eq.~\ref{eq:rom} strictly limits the possible EE movements to assure safety.
% Fig.~\ref{fig:mip_a1} shows the essential decision variables optimized in \ref{ocp:mip_general}. 
Due to the convex nature of this MICP, we can determine the feasibility of a possible symbolic transition by checking if the optimal solution exists.
% \begin{figure}
%     \centering
%     \includegraphics[width=0.75\linewidth]{mip_a1.pdf}
%     \caption{Caption}
%     \label{fig:mip_a1}
% \end{figure}


%%%% Fig4. System overview %%%%
\begin{figure*}[t]
    \centering
    \includegraphics[width=\linewidth]{combined_offline_online_v3.pdf}
    \caption{System overview for offline synthesis and online execution. During the offline synthesis (black dashed arrows), an initial set of locomotion gaits is given and symbolic skills are iteratively generated by solving MICP. When the task specifications are unrealizable, a symbolic repair is triggered to seek missing skills. During the online execution (black solid arrows), MICP is solved again taking online terrain segments and the symbolic state only advances when a solution is found. Online repair is triggered (blue solid arrow) when a solving failure or an unseen terrain is detected. 
    % both abstracted and real-world terrain information is fed into the automaton.
    }
    \label{fig:combined_offline_online}
\end{figure*}
%%%%%%%%%%%%%%%%%%%%%%%%%%%%%%%%%%%

% \subsection{GR(1) Specifications given Robot Skills}\label{sec:spec_gen}
\subsection{Task Specification Encoding}\label{sec:spec_gen}

After acquiring a set of feasible robot skills, 
we encode the skills into the specification $\spec$ 
as part of the environment safety assumptions $\envsafety$ 
and system safety guarantees $\syssafety$.
%
% We divide the safety assumptions $\envsafety$ into mutable assumptions $\envsafetynothard$,
% which describe the postconditions of the skills,
% and hard assumptions $\envsafetyhard$ that cannot be modified, 
% i.e.
% $\envsafety = \envsafetynothard \land \envsafetyhard$.
% Similarly, we divide the safety guarantees $\syssafety$ into mutable guarantees $\syssafetynothard$,
% which represent the preconditions of the skills,
% and hard guarantees $\syssafetyhard$ that cannot be modified
% ($\syssafety = \syssafetynothard \land \syssafetyhard$).
%
% To translate from states to formulas, 
We define a function $\formplain$ that takes as input
two subsets $\statevar, \statevar' \subseteq \ap$,
where $\statevar \subseteq \statevar'$,
and generates a Boolean formula
$\form{\statevar'}{\statevar} \coloneqq \bigwedge_{\var\in\statevar}\var \land \bigwedge_{\var \in \statevar'\setminus\statevar} \neg\var$.
For any subset $\statevar \subseteq \ap$,
we use $\bigcirc\statevar$ to denote the set $\{\bigcirc \pi \mid \pi \in \statevar\}$.

% We encode the postconditions of the skills into the environment safety assumptions $\envsafety$ of the specification $\spec$.



% each skill $\texttt{skill} \in \mathcal{A}$ comes with a precondition $\sum_{\texttt{skill}}^{\rm{pre}} \in 2^{\mathcal{I}_G}$ and a postcondition $\sum_{\texttt{skill}}^{\rm{post}} \in 2^{\mathcal{I}_G}$ similar to \cite{pacheck2022physically}. In Fig.~\ref{fig:preliminary_nine_squares}(b) and (c), it shows a skill $\texttt{skill0}$ with precondition $\sum_{\texttt{skill0}}^{\rm{pre}} = \{(x=1,y=2,\rm cT = \rm Flat, \rm rT = \rm High)\}$ and postcondition $\sum_{\texttt{skill0}}^{\rm{post}} = \{(x=2,y=2)\}$. The skills are encoded in the GR(1) formula as part of $\phi^\textrm{t}_\textrm{e}$ and $\phi^\textrm{t}_\textrm{s}$, 
% which can be further divided into $\phi^{\textrm{t}}_{\textrm{e}} = \phi^{\textrm{t, mutable}}_{\textrm{e}} \wedge \phi^{\textrm{t, hard}}_{\textrm{e}}$ and $\phi^{\textrm{t}}_{\textrm{s}} = \phi^{\textrm{t, mutable}}_{\textrm{s}} \wedge \phi^{\textrm{t, hard}}_{\textrm{s}}$. $\phi^{\textrm{t, mutable}}_{\textrm{e}}$ and $\phi^{\textrm{t, mutable}}_{\textrm{s}}$ capture the postconditions and preconditions that the repair process might modify, respectively. $\phi^{\textrm{t, hard}}_{\textrm{e}}$ and $\phi^{\textrm{t, hard}}_{\textrm{s}}$ encode the hard constraints specified by the user remain unchanged and always hold.   

% \subsubsection{Postconditions}
The skill assumptions $\envsafetynothard$ encode the postconditions of the skills. 
% As shown in Eq.~\ref{eq:env_trans}, 
The assumptions ensure that after the skill execution, 
at least one postcondition of the skill must hold:
% s in the next timestep:

\begin{equation}\label{eq:env_trans}
\begin{aligned}
    \envsafetynothard \coloneqq \bigwedge_{\skill \in \mathcal{O}} \big( \skill 
    % \land 
    % \form{\inpcen\cup\inpterrain}{\skillpre} \\ 
    \to
    \bigvee_{\statevar\in\skillpost} \form{\bigcirc\inpcen}{\bigcirc\statevar} \big)
    % \bigvee_{\sigma_\textrm{pre} \in \sum_{\texttt{skill}}^\textrm{pre}} \varphi_{\sigma_{\textrm{pre}}}) \rightarrow \\
    % &\bigcirc \bigvee_{\sigma_\textrm{post} \in \sum_{\texttt{skill}}^\textrm{post}} \varphi_{\sigma_\textrm{post}}\bigg)
\end{aligned}
\end{equation}
%
In Example~\ref{fig:preliminary_nine_squares},
% $\skillprecustom{\skillcustom{0}}$,
the postcondition of skill $\skillcustom{0}$
is defined as:
% encoded asin Eq.~\ref{eq:skill_0_post}:
% Therefore, as an example, the postcondition specification for $\texttt{skill0}$ is defined as:
\begin{equation}\label{eq:skill_0_post}
\begin{aligned}
     \square (\skillcustom{0} 
     % \land x_1 \land y_1 \land \tvar{x_1}{y_0} \land \tvar{x_1}{y_1} \land \tvar{x_1}{y_2}
     % \land \bigwedge_{\var\in\inpcen\cup\inpterrain\setminus\skillprecustom{\skillcustom{0}}} \neg \var \\
    \to &\bigcirc x_0 \land \bigcirc y_1 \land \neg \dots)
    % &\neg \bigcirc x_1 \land 
    % \neg \bigcirc x_2 \land 
    % \neg \bigcirc y_0 \land 
    % \neg \bigcirc y_2 \big)
     % \bigwedge_{\var\in\inpcen\setminus\skillpostcustom{\skillcustom{0}}} \neg \bigcirc \var \big)
     % x=1 \wedge y=2 \wedge \rm cT = \rm Flat \wedge \rm rT = \rm High \big) \rightarrow \\
     % &\bigcirc \big(x=2 \wedge y=2 \big) \big)
\end{aligned}
\end{equation}
where we omit the centroidal inputs that are False.

% \subsubsection{Preconditions}
The skill guarantees $\syssafetynothard$ encode the preconditions of the skills.
The guarantees impose restrictions on when a skill can be executed.
% As shown in Eq.~\ref{eq:sys_trans},
The guarantees only allow the robot to execute a skill 
if one of the skill's preconditions holds:
% The preconditions encode constraints in $\varphi_{\textrm s}^{\textrm{t}, \textrm{mutable}}$. It restricts the conditions that a specific skill can be chosen. Therefore, it allows the overall system to choose a skill only when a precondition is satisfied, but does not require the system to do so. 

\begin{equation}\label{eq:sys_trans}
\begin{aligned}
    \syssafetynothard \coloneqq 
    \bigwedge_{\skill \in \mathcal{O}} \bigg(
    \neg \big(
    \bigvee_{\statevar\in\skillprecustom{\skill}}
    \form{\bigcirc\inpcen\cup\bigcirc\inpterrain}{\bigcirc\statevar}
    \big)
    % \\
    % \form{\bigcirc(\inpcen\cup\inpterrain)}{\bigcirc\skillprecustom{\skill}}\\
    \to \neg \bigcirc \skill
    \bigg)
    % \\
    % \varphi_{\textrm s}^{\textrm{t}, \textrm{mutable}} = \bigwedge_{\pi_{\texttt{skill}} \in \mathcal{O}} \bigg( & \neg (\bigcirc \bigvee_{\sigma_\textrm{pre} \in \sum_{\texttt{skill}}^\textrm{pre}} \varphi_{\sigma_{\textrm{pre}}}) \rightarrow 
    % &\neg \bigcirc \pi_{\texttt{skill}}\bigg)
\end{aligned}
\end{equation}
%
In Example~\ref{example:1}, the precondition of skill $\skillcustom{0}$ is defined as:
\begin{equation}
\begin{aligned}
     \square \big( \neg 
     (
     \bigcirc x_1 \land 
     \bigcirc y_1 \land 
     \bigcirc \tvar{x_1}{y_0} \land 
     \bigcirc \tvar{x_1}{y_1} \land 
     \bigcirc \tvar{x_1}{y_2}\land \\
     \neg \dots)
     % \neg \bigcirc x_0 \land
     % \neg \bigcirc x_2 \land
     % \neg \bigcirc y_0 \land
     % \neg \bigcirc y_2 \land \\
     % \bigwedge_{i\in\{0,2\}, j\in\{0,1,2\}}
     % \neg \bigcirc \tvar{x_i}{y_j}\big)
     % \bigwedge_{\var\in\inpcen\cup\inpterrain\setminus\skillprecustom{\skillcustom{0}}} \neg \bigcirc \var
     % \big)
     % \\
     \to \neg \bigcirc \skillcustom{0} \big)
     % \\
     % (\bigcirc (x=1 \wedge y=2 \wedge \rm cT = \rm Flat \wedge \rm rT = \rm High)) \rightarrow \\
     % \neg \bigcirc \pi_\texttt{skill0}\big)
\end{aligned}
\end{equation}
where we omit the centroidal and terrain inputs that are False.
% include the keyframe transition action $\mathcal{K}$ and gait primitives $\mathcal{G}$. We further divide $\mathcal{G}$ into $\mathcal{G} = \{\rm{gaitType}, \rm{gaitSpeed}\}$. We incorporate extra input propositions $\mathcal{E}$ to represent the terrain information, which shows up in the preconditions as well.
% \subsubsection{Hard Constraints}\label{sec:hard_constraints}

The hard constraints ($\envsafetyhard$, $\syssafetyhard$) are constraints that a repair cannot modify 
(see Sec.~\ref{sec:symbolic_repair}).
To strictly guarantee $\syssafetyhard$,
we only allow the robot to execute one skill at a time:
    $\square\big(\neg (\bigcirc\skill \land \bigcirc\skill')\big)$, for any two different skills $\skill, \skill' \in \out$. Given that, 
    %such that $\skill \neq \skill'$.
we encode the following hard assumptions $\envsafetyhard$:
% as hard constraints in $\envsafetyhard$ and $\syssafetyhard$:
% The hard assumptions $\envsafetyhard$ states that
\begin{enumerate}[label=\arabic*),ref=\arabic*]
    \item The uncontrollable terrain and request inputs 
    % $\inpterrain \cup \inprequest$
    % which include the terrain and request inputs ($\inpuser = \inpterrain\cup\inprequest$), 
    cannot change during execution:
    % $\forall \var \in \inpuser$,
    $\square (\var \leftrightarrow \bigcirc \var)$,
    $\forall \var \in \inpterrain \cup \inprequest$.
    \item The centroidal inputs $\inpcen$ remain unchanged if no skill is executed: 
    $\square\big(\bigwedge_{\skill\in\out} \neg \skill \to \bigwedge_{\var\in\inpcen} (\var \leftrightarrow \bigcirc \var)\big)$.
    \item \label{assumption:poss_terrain}
    The terrain state must come from the set of possible terrain states $\inpstateterrainposs \subseteq 2^{\inpterrain}$ provided by the user:
    $\envterrain \coloneqq \square (
    \bigvee_{\statevar\in\inpstateterrainposs}
    \form{\inpterrain}{\statevar}
    )$.
\end{enumerate}

% Certain constraints also need to be incorporated as hard constraints. $\varphi_{\textrm e}^{\textrm{t}, \textrm{hard}}$ and $\varphi_{\textrm s}^{\textrm{t}, \textrm{hard}}$. For instance, specifications such as enforcing movements among adjacent grids are encoded as hard constraints. Other constraints such as the mutual exclusion of different skills are also added here. 

% For the nine-square example, we 

% Due to the space limit, we only highlight the specifications that build connections between the environment and the robot action. We construct the following specifications based on locomotion heuristics:

% Hard specifications. React to terrain abstraction. Mutual exclusion of each skill and each grounded state.



% \subsection{Offline Synthesis with Feasibility Check}
% Fig.~\ref{fig:combined_offline_online} illustrates the system overview for both offline synthesis and online execution. This section focuses on the offline synthesis process. Using prior knowledge provided by humans, we first define a set of motion primitives, such as a one-second trotting gait, and a set of potential terrains described at the symbolic level. Based on the abstracted robot and terrain states defined in Sec.\ref{sec:preliminaries_abstraction}, we iterate over the grid locations and solve the MICP to check if a solution exists when the desired goal is the adjacent grid. If a solution is found, the corresponding skill is added to the specification. Once the hard constraints and physically feasible skills are incorporated into the specification, we assess its realizability. If it is realizable, a strategy is generated for online execution. If not, an offline repair is introduced, as discussed in the next section.

\subsection{Symbolic Repair and Gait-Free MICP}\label{sec:symbolic_repair}
The specification $\spec$ generated in Sec.~\ref{sec:spec_gen} may be unrealizable
if the physical feasibility checking in Sec.~\ref{sec:motion_primitive} does not generate enough skills for the robot to reach all goal locations over all possible terrain combinations.
%
To adress this issue, we leverage the ideas in~\cite{pacheck2022physically} 
% to ensure that the generated specification $\spec$
% (Sec.~\ref{sec:spec_gen}) is realizable~\cite{pacheck2022physically}.
% The specification $\spec$ can be unrealizable if the physical feasibility checking (Sec.~\ref{sec:motion_primitive}) does not generate enough skills for the robot to traverse to all goal locations over all possible terrain combinations.
% The repair 
to
generate
new skill suggestions necessary for the robot to accomplish the task.
%over all possible terrain environments.
Different from~\cite{pacheck2022physically}, we assume that skills do not contain intermediate states,
and the skill preconditions
are over both the controllable centroidal inputs $\inpcen$ and uncontrollable terrain inputs $\inpterrain$, 
while the postconditions remain to contain 
controllable
centroidal inputs $\inpcen$ only.

% We extend repair to handle skills whose preconditions are over both the controllable centroidal inputs $\inpcen$ and terrain inputs $\inpterrain$ while the postconditions still only contain controllable inputs $\inpcen$.

For any subset $\statevar\subseteq \ap$, we say that $\statevar$ is a winning state
if the robot can guarantee to satisfy the task from $\statevar$
despite adversarial environment behaviors.
The repair process systematically expands the set of winning states until
it intersects with all liveness goals and contains all initial states.
To this end, the repair alternatively modifies the preconditions and postconditions of existing skills.
When modifying preconditions, the repair randomly selects a skill whose postcondition is winning, and modifies its precondition to a non-winning state, thus making the new precondition possible to win.
When modifying postconditions, the repair randomly selects a skill whose preconditions and postconditions are both non-winning, and modifies its postcondition to a winning state, thus making the original precondition possible to win.
%
Fig~\ref{fig:preliminary_nine_squares}c illustrates the use of repair in Example~\ref{example:1}. 
To create a skill that reaches the requested grid $(x_0, y_0)$, the
repair selects the skill transition that moves the robot from the grid $(x_1, y_0)$ to $(x_2, y_0)$,
and modifies the postcondition to reach $(x_0, y_0)$ instead.

We implement the suggested skills by reconfiguring the gait in terms of contact sequence to generate new motion primitives. This can be achieved either by leveraging human knowledge to introduce new candidate gait parameters or by solving a gait-free version of Eq.~\ref{ocp:mip_general}. The gait-free MICP retains all continuous decision variables and constraints from the gait-fixed MICP but modifies the contact state-dependent constraints, as the contact sequence and timing are no longer fixed. Instead of using $\v H_{r,k,j}$, we introduce a different set of binary variables $\v H_{r,m,j}$ for defining the contact state for each foot at each time step explicitly similar to \cite{valenzuela2016mixed}. $r$ and $j$ still represent the $r^{\rm th}$ convex region and $j^{\rm th}$ foot, respectively, while $m \in \range{M}$ denotes the $m^{\rm th}$ time step, evenly discretized with $\Delta t_m$ over the entire $M$ time steps. Therefore, the gait-free MICP involves introducing additional binary variables to determine contact states, which, however, requires a longer computation time.
If any suggested skill is not feasible, 
we add the symbolic transition as a hard constraint on repair,
and repeat the repair process.

% % \subsection{Outline}
% \begin{enumerate}
%     \item We define a state $\statevar \subseteq \ap$ to be winning,
%     and thus $\statevar$ is a winning state,
%     if the robot can guarantee to satisfy the task goals from $\statevar$ despite adversarial environment behaviors
%     % \item We define the set of winning states as $\winning$
%     \item Repair systematically expands the set of winning states by alternatively modifying the pre and postconditions of existing skills.
%     \item When modifying postconditions, repair randomly selects a transition whose pre and post-conditions are both not winning, and modifies the postcondition to a winning state, thus making the precondition also possible to win.
%     \item When modifying preconditions, repair randomly selects a transition whose postcondition is winning, and modifies the precondition to a not-winning state, thus making the new precondition more likely to win.
%     \item Fig~\ref{fig:preliminary_nine_squares}c illustrates the use of repair in Example~\ref{example:1}. 
%     Repair selects the skill that moves the robot from the grid $(x_1, y_0)$ to $(x_2, y_0)$,
%     and modifies the postcondition to be $\{x_0, y_0\}$.
%     \item 
%     In Example~\ref{example:1}, the robot does not have enough skills to reach the requested goal grid $(x_0, y_0)$, as shown in 
%     Fig~\ref{fig:preliminary_nine_squares}c. 
%     Repair randomly selects the skill $\skillcustom{1}$ whose precondition is in $(x_1, y_0)$ and postcondition is in $(x_2, y_0)$,
%     and modifies 
%     % - states from which the robot can guarantee to satisfy the task goals.
%     % This extension allows the skills 
%     % \begin{enumerate*}[label=(\roman*),ref=\roman*]
%         % \item skills do not contain intermediate states,
%         % \item we only add new skills instead of modifying existing skills, and
%         % \item 
%     % \end{enumerate*}
    
%     % the specification $\spec$ generated in Sec.~\ref{sec:reactive_gait_synthesis} becomes unrealizable.
%     % \item We utilize a symbolic repair tool to generate new robot skills that make $\spec$ realizable~\cite{pacheck2022physically}.
    
    
% \end{enumerate}
% When the specifications above are not realizable, we perform a repairing process (green arrows in Fig.~\ref{fig:combined_offline_online}). In contrast to \cite{pacheck2022physically}, we assume: (i) no intermediate state is considered for existing skills; (ii) the existing skills remain unchanged and we only add new skills; (iii) the new skill's precondition and postcondition are adjacent to each other; (iv) we incorporate extra input propositions $\mathcal{E}$ to represent the terrain information, which shows up in the preconditions as well. During the repair, human is obligated to provide new motion primitives to restore the missing skills found by the repair. The human can forfeit the current symbolic suggestion and move on to the next one. 

% \subsection{New Skills Generation}
% Preconditions and postconditions modification.
% \begin{itemize}
%     \item Reactive inputs pre-processing.
%     \item Terrain inputs.
% \end{itemize}

% \subsection{Gait Optimization}
% (TODO) To potentially add the new skills suggested by the symbolic 

% \subsection{Semi-Automatic Repair: Human in the Loop}
% The human is in charge of collecting more locomotion primitives to allow for more feasible transitions.

\section{Online Execution}\label{sec:online_execution}
The online execution module takes as input
a terrain state $\inpstateterrain \subseteq \inpterrain$
from the perception module, 
a request state $\inpstaterequest \subseteq \inprequest$
from the global planner,
and a strategy $\strategy$
from the offline synthesis module (Sec.~\ref{sec:reactive_gait_synthesis}).
As shown in Fig.~\ref{fig:combined_offline_online},
if the robot encounters an unseen terrain state, 
% is unexpected and thus not considered by the offline synthesis module,
we leverage symbolic repair to create new skills that handle the unexpected terrain state online 
% $\inpstateterrain$ 
% and re-synthesize a strategy
(Sec.~\ref{sec:online_repair}).
This module then executes the strategy automaton $\strategy$
by solving a MICP problem again 
with more detailed and accurate terrain information 
to generate a nominal trajectory for each transition (Sec.~\ref{sec:automaton_online_evaluation}). This nominal trajectory will be further tracked by a tracking controller in real time
(Sec.~\ref{sec:tracking_control}).
After reaching the request state $\inpstaterequest$,
the robot resets the strategy with new terrain, requests states,
and repeats the execution
until reaching the final goal state.
% We formulate a MICP problem using 





% During the online execution, we consider the mismatch between the abstracted and continuous information including both the robot and environmental states. 
% The online execution is divided into an online strategy execution process and a tracking control module.

\textbf{Given: 1. terrain state, 2. request state, 3. disallowed transitions}
\subsection{Online Repair}\label{sec:online_repair}
During online execution, 
the robot may encounter unexpected terrain states
that the initial model of the workspace structure does not take into consideration.
We leverage an online repair tool~\cite{meng2024automated} 
to detect and recover from such unexpected terrain states.

The online repair tool first generates a monitor that detects violations of the environment assumptions $\envsafety$ during online execution.
%
Note that $\envsafety$ includes the assumption
$\envterrain = $ $
    \square (
    \bigvee_{\statevar\in\inpstateterrainposs}
    \form{\inpterrain}{\statevar})$,
which assumes that 
at any time, 
at least one terrain state from the set of possible terrain states $\inpstateterrainposs$ must hold
% the current terrain state $\inpstateterrain$
% must always come from
% $\inpstateterrainposs$, 
% the set of possible terrain states provided by the user
(see Sec.~\ref{sec:spec_gen}).
% Suppose the current terrain state $\inpstateterrain$ is not in $\inpstateterrainposs$, the set of possible terrain states provided by the user.
During online execution,
if the current terrain state $\inpstateterrain \not\in \inpstateterrainposs$,
the monitor detects that $\inpstateterrain$
violates the assumption $\envterrain$.
We then relax the violated assumption $\envterrain$ to admit the new terrain state $\inpstateterrain$,
replacing $\envterrain$ with 
$\square (
    \bigvee_{\statevar\in\inpstateterrainposs\cup\{\inpstateterrain\}}
    \form{\inpterrain}{\statevar})$.
After assumption relaxation,
the revised specification becomes unrealizable
as the robot lacks the required skills to handle the new terrain state $\inpstateterrain$.
We utilize symbolic repair 
(Sec.~\ref{sec:symbolic_repair})
to find a set of necessary skills for the new terrain state.






% the prior knowledge of the workspace structure does not take into account.
% \begin{enumerate}
%     \item Encounter unexpected terrain states during execution
%     \item Generate a monitor to check assumption violations
%     \item The monitor detects that the new terrain state violates the environment assumption which states that the current terrain state must belong to the set of possible terrain states.
%     \item we relax the assumption to include the new terrain state
%     \item we perform repair obtain new skills that can handle the new terrain state, and resume the execution
% \end{enumerate}

\subsection{Strategy Execution}\label{sec:automaton_online_evaluation}
The red path in Fig.~\ref{fig:combined_offline_online}(c) represents the automaton online evaluation process. Before transitioning to the new symbolic state, an online MICP is solved. The automaton advances only if the MICP successfully finds a solution.
% Algorithm \ref{code:online_automaton} shows the complete steps of performing online evaluation. More specifically, we consider multiple.
% \subsubsection{Online MICP}
Slightly different from the MICP formulation in Eq.~\ref{ocp:mip_general} during the offline phase, the online MICP is executed given accurate terrain information  from a terrain segmentation module. Since a strategy is already generated offline, the online MICP explicitly knows the chosen gait and the ideal initial and final robot states given by the robot and terrain abstraction.

% \subsubsection{Robot Pose Re-Targeting}
% Before running TO, the final state will be re-evaluated by solving a kinematic feasibility problem similar to Eq.~(1). The dynamics equation (\ref{eq:dynamics}) is replaced by enforcing the center of mass (CoM) to lie inside a convex support polygon to ensure static stability. In addition, a hard constraint is added to ensure that the modified robot pose stays within a certain threshold compared with the original desired pose. The cost function is simplified to stay as close as possible to the original desired pose. The higher-order terms of body position, orientation, and end-effectors are excluded from the decision variables. All the other constraints remain the same. 

% \subsubsection{Failure Handling}
% Replan process.

\subsection{Tracking Control}\label{sec:tracking_control}
\subsubsection{Nonlinear Model Predictive Control}\label{subsec:nmpc}
The nonlinear model predictive controller (NMPC) operates independently from the symbolic planning module to track the reference trajectory generated by the online MICP, utilizing a more accurate dynamics model. In this work, we use a NMPC considering the centroidal dynamics and kinematics of the robot similar to \cite{dai2014whole,sleiman2021unified}. The optimization problem is formulated as a Sequential Quadratic Program (SQP) and solve through
%by linearizing the nonlinear constraints present in the formulation such as the dynamics. 
the OCS2 library \cite{OCS2}. We employ a time horizon of 1 second with 100 timesteps.

% \subsection{Coordination between Automaton and NMPC }
% \begin{figure}
%     \centering \includegraphics[width=\linewidth]{multi_threading_v2.pdf}
%     \caption{MPC and Automaton Coordination: The trajectory windows are color-coded as follows: green represents trajectories generated by NMPC with a shorter time horizon but faster frequency; grey indicates trajectories generated by MICP with a longer time horizon, spanning from one symbolic node to another under ideal initial conditions; red denotes trajectories generated by MICP after replanning, triggered when the predicted state from MPC deviates from the expected goal condition.}
%     \label{fig:micp_mpc_automaton}
% \end{figure}
% Since a tracking error is hard to avoid, a harmonic coordination between these two threads is necessary. The NMPC is running along with the automaton feasibility checker as shown in Fig.~\ref{fig:micp_mpc_automaton}. The overall process can be summarized in the following steps:

% \subsubsection{Automaton Execution} The automaton thread initiates the process by solving MICP and generating reference trajectories as decribed in \ref{subsec: automaton_online_evaluation}.

% \subsubsection{MPC Tracking}
% Once the initial reference trajectory is generated, the MPC thread begins tracking it using a more accurate dynamics model as described in \ref{subsec:nmpc}.

% \subsubsection{Error Handling and Replanning}
% If the MPC thread predicts a state that reaches the final state but deviates beyond the specified tolerance, the automaton thread resets to this new initial state and reruns the MICP to adjust the plan. If no significant deviation occurs, the automaton thread continues rolling out, using the desired state until the final waypoint is reached.

\subsubsection{Whole-Body Control}
A whole-body controller (WBC) is deployed to track the NMPC trajectory plan by solving a hierarchical quadratic program (QP) \cite{herzog2016momentum} at 500 Hz. The state estimator fuses IMU data, joint encoders, and motion capture inputs to provide accurate body position information.
% \subsection{Integration to Higher-Level Command}
% Reset process.

\section{Results}\label{sec:experiments}
We present case studies to demonstrate the synthesis and repair processes. Examples of maneuvering over various environments highlight the framework's efficacy, scalability, and generalizability. Hardware experiments are demonstrated to show the real-world deployment.
\subsection{Robot Setup}
As mentioned earlier, we verify our algorithms on two separate robot platforms, Unitree Go2 and SkyMul Chotu. The parameters for each robot are defined in Table~\ref{tab:robot_parameter} and used for formulating our optimization problems, including the robot mass $m$, torque limit $\vg \tau_{\text{max}}$ and the reference joint position $\v q_j^{\text{ref}}$ for a single leg with three joints, the reference foot position relative to the base frame $^\mathcal{B} \v p_j^{\text{ref}}$ and the maximum deviation of the foot position $\v p_j^{\text{max}}$ for the front left leg ($j=0$), and the moment of inertia of the approximated single rigid body $\v I$.
\begin{table}[t]
    \caption{Robot Parameters}
    \centering
     \begin{tabular}{c | c | c}
         \hline
         Parameter & Unitree Go2 & SkyMul Chotu\\\hline
         $m (kg)$ & 15.071 & 19.978\\
         $\vg \tau_{\text{max}} (N \cdot m)$ & (23.5, 23.5, 45.4) & (23.5, 23.5, 33.5)\\
         $\v q_j^{\text{ref}}(rad)$ & (0.0, 0.72, -1.44) & (0.0, 0.72, -1.44)\\
         $^\mathcal{B} \v p_j^{\text{ref}}(m)$ & (0.1805, 0.1308, -0.29) & (0.2118, 0.210, -0.30)\\
         $\v p_j^{\text{max}}(m)$ & (0.15, 0.1, 0.15) & (0.12, 0.08, 0.15)\\
         $\v I (kg \cdot m^2)$ & diag(0.152, 0.369, 0.388) & diag(0.396, 0.915, 1.107)\\
         \hline
    \end{tabular}
    \label{tab:robot_parameter}
\end{table}

\subsection{Maneuvering Showcase}
To demonstrate the generalizability of the proposed framework, we examine two scenarios with varying terrain types, numbers of potential safe stepping regions, and robot platforms. The \textit{obstacle} is also explicitly captured as a terrain type and applied to all scenarios. To avoid transitioning to an obstacle grid, we add specifications for obstacle avoidance as hard constraints, and assume the requested waypoint is not an obstacle grid. The abstracted grid size can be extended to higher resolution and for all examples we study $3 \times 3$ and $5 \times 5$ grid size separately. The $5 \times 5$ grid could have longer planning horizon and more decisions to make but also increase the problem complexity. We use a 0.8 m cell size for the unstructured terrain and 0.6 m for the rebar scenario, but this can be adjusted for practical applications.

\subsubsection{Unstructured Terrain}
% \begin{figure}
%     \centering
%     \includegraphics[width=\linewidth]{online_terrains.pdf}
%     \caption{(Will be updated to have rebar scenario.) Online scenarios with random sparsity of stepping stone and obstacles.}
%     \label{fig:online_terrains}
% \end{figure}
Unlike Example \ref{example:1}, to further study the scalability of the proposed framework, we abstract eight terrain types: \textit{flat terrain}, \textit{high terrain}, \textit{low terrain}, \textit{dense stone}, \textit{sparse stone}, \textit{gap}, \textit{high gap} and \textit{low gap}. The classification criterion is heuristically determined based on polygon height, the number of polygons, and their area relative to the abstraction cell.  We define two example specifications that consider 4 and 8 terrain types selected from the aforementioned terrain types, and create corresponding scenarios as shown in Fig.~\ref{fig:online_terrains}.
% \begin{figure}[h]
%     \centering
%     \includegraphics[width=\linewidth]{5x5_unstructured_demonstration.pdf}
%     \caption{Local environment abstraction and maneuvering with 5x5 (upper) and 3x3 (lower) grid size in separate unstructured terrain scenarios.}
%     \label{fig:multiple_sizes}
% \end{figure}

\subsubsection{Rebar}
Inspired by the labor-intensive rebar tying task on construction sites, we explore a second application where the quadrupedal robot navigates a rebar mat. Each rebar is modeled as a rectangular polygon with a 3 cm width. Unlike the stepping stone scenario, this setup offers more potential stepping regions due to the higher rebar density. Observing that the robot's traversability depends on rebar sparsity in different directions—whether aligned with its facing direction or not given a tolerance—we abstract and classify rebar types based on a combination of horizontal sparsity (perpendicular to the robot's facing direction) and vertical sparsity (parallel to it). Each direction's sparsity is categorized as \textit{none}, \textit{single}, \textit{dense}, \textit{sparse}, or \textit{extreme sparse}, determined heuristically by the number and relative spacing of the rebars inside the abstraction cell. To reduce problem complexity, some trivial combinations are removed such as \textit{none} or \textit{single} in both directions. As a result, we define two example specifications with 7 and 14 rebar terrain types and corresponding real scenarios. The 7-type scenario generates random rebar sparsity between 0.15 and 0.35 m, excluding the extreme sparse type (above 0.4 m). In contrast, the 14-type scenario includes rebar sparsity from 0.15 to 0.6 m, potentially covering more non-trivial combinations including the extreme sparse ones.
Fig.~\ref{fig:rebar-traversing} illustrates the rebar scenario.

\begin{figure}
    \centering
    \includegraphics[width=\linewidth]{unstructured_8_terrain.pdf}
    \caption{Unstructured terrain scenarios with random sparsity of stepping stone and obstacles. \todo make it more interesting (obstacle, randomness).}
    \label{fig:online_terrains}
\end{figure}


\begin{figure}
    \centering
    \includegraphics[width=\linewidth]{rebar_type_maneuver.pdf}
    \caption{Rebar scenarios.}
    \label{fig:rebar-traversing}
\end{figure}

The offline synthesis for both \textit{Unstructured Terrain} and \textit{Rebar} starts with a 2.0-second and a 3.0-second trotting gaits. 
% Note that different from the case study that mainly for demonstrating the process, we add this additional longer horizon trotting gait to allow more safe walking options from a practical standpoint. 
We generate a set of terrain states by systematically sweeping a local $3 \times 3$ or $5 \times 5$ grid across the entire terrain map using a predefined increment.
The request states are defined as the top row of the local grid map in the robot's facing direction.
The offline repair and feasibility checking time is reported in Table \ref{tab:offline_synthesis_time}. It is observed that a set of leaping gaits with a short aerial phase are generated by gait-free MICP for the newly found skills that involve transitioning from a lower terrain to a higher terrain, or cross a gap.
Table \ref{tab:offline_synthesis_time} shows the computational time for offline synthesis. 
% \todo A figure shows the repair time with regard to the number terrain states (highlight the complexity of repair). To further demonstrate the complexity of the repair process, we intentionally change the increment when sweeping through the terrain map, generating different numbers of terrain states for offline synthesis and repair. This approach is reasonable when global terrain knowledge is limited.

Fig.~\ref{fig:online_terrains} shows the visualization snapshots of the Go2 and Chotu traversing aforementioned terrains during online execution. The robot is tasked with reaching a global waypoint using a naive global planner to find the shortest path.
%selecting the closest grid in the 9-grid map as the local waypoint. 
Specifically, the local waypoint in the local grid map is selected by choosing the closest, non-obstacle grid toward the global waypoint.
The top sub-figure in Fig.~\ref{fig:online_terrains} shows a sequential maneuver through flat terrain, dense stepping stone, sparse stepping stone, and a gap. The bottom one shows the robot navigates through all 8 terrain types mentioned above by adaptively choosing the locomotion gait.
Table \ref{tab:micp_speed} shows the online computation speed for MICP.

\subsection{Online Repair Case Study}
As a case study to demonstrate the capability of handling unforeseen terrains and failures, we highlight the following run-time scenarios that trigger an online repair:
\subsubsection{Unforeseen Terrain States} With limited knowledge of the global terrain map, the terrain states provided during offline synthesis may be insufficient when encountering new terrain configurations. Fig.~\ref{fig:online_repair} illustrates this in an Unstructured - 4 types scenario with a $3 \times 3$ grid size, where terrain states involving gap terrain are excluded, treating gaps as an unseen terrain type during offline synthesis. When the robot reaches a new terrain state that includes a gap, an online repair is triggered for the current terrain and request states. A new skill with a leaping gait is generated to successfully cross the gap. The online repair takes \todo, where the pure symbolic repair part takes \todo and the gait-free MICP takes \todo to generate the new leaping gait. 
\begin{figure}
    \centering
    \includegraphics[width=\linewidth]{online_repair_new_gap.pdf}
    \caption{Caption}
    \label{fig:online_repair}
\end{figure}

\subsubsection{Solving Failure} Another common runtime failure arises from MICP solving failures. Due to discrepancies between the predefined polygons used in offline synthesis and the actual online terrains—such as variations in shape and size—a skill deemed feasible offline may become infeasible when executed online, even if classified under the same symbolic transition. The detours shown in Fig.~\ref{fig:rebar-traversing} illustrate the online repair and re-synthesis process following solving failures in both the 7-type and 14-type cases. This issue is observed more frequently in the rebar scenario compared to unstructured terrain. The online repair takes \todo and \todo, which is significantly shorter than the previous case due to the reuse of existing skills.



% we preclude the evaluation for terrain (d). An online scenario is then created as shown in Fig.~\ref{fig:online_repair} with randomized sparsities ranging from 0 to 0.15 m. In Fig.~\ref{fig:online_repair}, the robot is tasked with reaching a global waypoint using a naive global planner to find the shortest path.
% %selecting the closest grid in the 9-grid map as the local waypoint. 
% Then the local waypoint in the 9-grid map is selected by choosing the closest grid. 
% The robot starts to navigate until it encounters an unforeseen terrain case (e) after abstracting the perceived local terrain shapes. Grids $(x0,y0)$ and $(x2,y0)$ are unreachable as shown in Fig.~\ref{fig:online_repair}(e1) since the sideway walking capability between sparse terrains was not captured. Upon detecting such a violation of the environment specification, the online repair is triggered, generating a new dynamically feasible skill for transitions $(x1,y0) \rightarrow (x0,y0)$ and $(x2,y1) \rightarrow (x2,y0)$ in \todo seconds shown in Fig.~\ref{fig:online_repair}(e2). The robot then successfully traverses terrain case (e) and can handle this terrain when encountering it again.

% (Done) Show the online repair result for a certain type of terrain. (Online terrain snapshot $\rightarrow$ polygons considered in the 9-grid $\rightarrow$ abstraction $\rightarrow$ existing skills $\rightarrow$ repaired skills)
% \begin{figure}
%     \centering
%     \includegraphics[width=1.0\linewidth]{online_repair.pdf}
%     \caption{Online repair results after witnessing unforeseen terrains.}
%     \label{fig:online_repair}
% \end{figure}

\begin{table*}[t]
    \caption{Offline Synthesis and Repair Time}
    \centering
     \begin{tabular}{| c | c | c | c | c | c | c}
         \hline
         Scenario & Grid Size & \begin{tabular}{@{}c@{}}Terrain States $\times$ Request States \\ (Success/Total)\end{tabular} & \begin{tabular}{@{}c@{}}Number of Skills \\ (Original/New/Total)\end{tabular} & Symbolic Repair Time (s) & \begin{tabular}{@{}c@{}}Feasibility Check Time \\ (Gait-Fixed/Gait-Free)\end{tabular}\\\hline
         Unstructured - 4 & $3 \times 3$ & -/48 & 13/\textcolor{red}{10}/64 & 5.99 &  14.77/- \\
         & $5 \times 5$ & -/365 & 18/\textcolor{red}{18}/64 & 225.33 &  25.11/- \\\hline
         Unstructured - 8 & $3 \times 3$ & -/75 & 11/\textcolor{red}{15}/256 & 9.26 & 15.87/- \\
         & $5 \times 5$ & -/500 & 26/-/256 & - & 40.53/-\\\hline
         Rebar - 7& $3 \times 3$ & 
         % -/258 & 47/-/52 & - &  116.69/- 
         -/93 & 32/\textcolor{red}{7}/196 & 15.27 &  86.36/- 
         \\
         & $5 \times 5$ & -/625 & 50/\textcolor{red}{14}/196 & 557.16 &  129.82/- \\\hline
         Rebar - 14 & $3 \times 3$ & -/96 & 38/\textcolor{red}{15}/784 & 19.46 & 114.09/-\\
         & $5 \times 5$ & -/745 & 86/-/784 & - & 295.74/- \\\hline
    \end{tabular}
    \label{tab:offline_synthesis_time}
\end{table*}

\begin{table*}[t]
    \caption{Online Gait-Fixed MICP Solving Time}
    \centering
     \begin{tabular}{|c | c | c | c | c | c | c}
         \hline
         Scenario & Collision Avoidance &
         Binary Variables &
         Continuous Variables & 
         Number of Polygons &
         Solve Time\\\hline
         Unstructured - 4 & No & 52 & 6583.5 & 6.5 & 0.37\\\hline
         Unstructured - 8 & Yes & 256 & 5166 & 2 & 2.65 \\
         & No & 69 & 7938 & 7.5 & 0.49\\\hline
         Rebar - 7& No & 70 & 7938 & 7.875 & 1.11\\
         \hline
         Rebar - 14 & No & 58 & 8142.75 & 6.25 & 1.09\\\hline
    \end{tabular}
    \label{tab:micp_speed}
\end{table*}

\subsection{Hardware Demonstration}
We demonstrate hardware implementations in a real-world rebar scenario. To ensure stable stepping and provide additional tolerance on the stepping error, each foot is equipped with a x-shape rubber. Fig.~\ref{fig:hardware_demo} shows a successful maneuver of Chotu on a rebar mat with non-homogeneous layout. Similar to the simulation examples, a global waypoint is sent and the robot continuously navigates in every $3 \times 3$ local grid map until the global goal is reached.
\todo Show the online repair after solving failure? Put an obstacle to trigger the online repair.

\begin{figure}
    \centering
    \includegraphics[width=\linewidth]{hardware_temp.pdf}
    \caption{Placeholder for hardware experiments.}
    \label{fig:hardware_demo}
\end{figure}

\section{Conclusion}
In this work, we presented an integrated planning framework for terrain-adaptive locomotion that combines reactive synthesis with MICP, enabling safe and reactive responses in dynamically changing environments. 
% To address unrealizable specifications arising from limited motion primitives, we introduced a symbolic repair approach that incorporates dynamic feasibility checks, automatically identifying missing transitions necessary for navigating adversarial terrains. In the online execution phase, we tackled the disparity between offline synthesis and real-world conditions by using an online MICP solver and an online symbolic repair process based on real-world terrain data. New terrain conditions encountered during execution. 
Our approach not only enhances motion feasibility and safety but also provides a scalable solution for legged robots maneuvering in complex and safety-critical environments. Future work will involve extensive hardware testing to validate the proposed framework in real-world scenarios. Additionally, we plan to explore specifications of greater complexity to further test the algorithmic scalability.


\bibliographystyle{IEEEtran}
\bibliography{references}

\end{document}
